%  LaTeX support: latex@mdpi.com
%  For support, please attach all files needed for compiling as well as the log file, and specify your operating system, LaTeX version, and LaTeX editor.

%=================================================================
\documentclass[bdcc,article,submit,pdftex,moreauthors]{Definitions/mdpi} 
%\documentclass[preprints,article,submit,pdftex,moreauthors]{Definitions/mdpi} 
% For posting an early version of this manuscript as a preprint, you may use "preprints" as the journal. Changing "submit" to "accept" before posting will remove line numbers.

% Below journals will use APA reference format:
% admsci, aieduc, behavsci, businesses, econometrics, economies, education, ejihpe, famsci, games, humans, ijcs, ijfs, journalmedia, jrfm, languages, psycholint, publications, tourismhosp, youth

% Below journals will use Chicago reference format:
% arts, genealogy, histories, humanities, jintelligence, laws, literature, religions, risks, socsci

%--------------------
% Class Options:
%--------------------
%----------
% journal
%----------
% Choose between the following MDPI journals:
% accountaudit, acoustics, actuators, addictions, adhesives, admsci, adolescents, aerobiology, aerospace, agriculture, agriengineering, agrochemicals, agronomy, ai, air, algorithms, allergies, alloys, amh, analytica, analytics, anatomia, anesthres, animals, antibiotics, antibodies, antioxidants, applbiosci, appliedchem, appliedmath, appliedphys, applmech, applmicrobiol, applnano, applsci, aquacj, architecture, arm, arthropoda, arts, asc, asi, astronomy, atmosphere, atoms, audiolres, automation, axioms, bacteria, batteries, bdcc, behavsci, beverages, biochem, bioengineering, biologics, biology, biomass, biomechanics, biomed, biomedicines, biomedinformatics, biomimetics, biomolecules, biophysica, biosensors, biosphere, biotech, birds, blockchains, bloods, blsf, brainsci, breath, buildings, businesses, cancers, carbon, cardiogenetics, catalysts, cells, ceramics, challenges, chemengineering, chemistry, chemosensors, chemproc, children, chips, cimb, civileng, cleantechnol, climate, clinbioenerg, clinpract, clockssleep, cmd, cmtr, coasts, coatings, colloids, colorants, commodities, complications, compounds, computation, computers, condensedmatter, conservation, constrmater, cosmetics, covid, crops, cryo, cryptography, crystals, csmf, ctn, curroncol, cyber, dairy, data, ddc, dentistry, dermato, dermatopathology, designs, devices, diabetology, diagnostics, dietetics, digital, disabilities, diseases, diversity, dna, drones, dynamics, earth, ebj, ecm, ecologies, econometrics, economies, education, eesp, ejihpe, electricity, electrochem, electronicmat, electronics, encyclopedia, endocrines, energies, eng, engproc, ent, entomology, entropy, environments, epidemiologia, epigenomes, esa, est, famsci, fermentation, fibers, fintech, fire, fishes, fluids, foods, forecasting, forensicsci, forests, fossstud, foundations, fractalfract, fuels, future, futureinternet, futureparasites, futurepharmacol, futurephys, futuretransp, galaxies, games, gases, gastroent, gastrointestdisord, gastronomy, gels, genealogy, genes, geographies, geohazards, geomatics, geometry, geosciences, geotechnics, geriatrics, glacies, grasses, greenhealth, gucdd, hardware, hazardousmatters, healthcare, hearts, hemato, hematolrep, heritage, higheredu, highthroughput, histories, horticulturae, hospitals, humanities, humans, hydrobiology, hydrogen, hydrology, hygiene, idr, iic, ijerph, ijfs, ijgi, ijmd, ijms, ijns, ijpb, ijt, ijtm, ijtpp, ime, immuno, informatics, information, infrastructures, inorganics, insects, instruments, inventions, iot, j, jal, jcdd, jcm, jcp, jcs, jcto, jdad, jdb, jeta, jfb, jfmk, jimaging, jintelligence, jlpea, jmahp, jmmp, jmms, jmp, jmse, jne, jnt, jof, joitmc, joma, jop, jor, journalmedia, jox, jpbi, jpm, jrfm, jsan, jtaer, jvd, jzbg, kidney, kidneydial, kinasesphosphatases, knowledge, labmed, laboratories, land, languages, laws, life, lights, limnolrev, lipidology, liquids, literature, livers, logics, logistics, lubricants, lymphatics, machines, macromol, magnetism, magnetochemistry, make, marinedrugs, materials, materproc, mathematics, mca, measurements, medicina, medicines, medsci, membranes, merits, metabolites, metals, meteorology, methane, metrics, metrology, micro, microarrays, microbiolres, microelectronics, micromachines, microorganisms, microplastics, microwave, minerals, mining, mmphys, modelling, molbank, molecules, mps, msf, mti, multimedia, muscles, nanoenergyadv, nanomanufacturing, nanomaterials, ncrna, ndt, network, neuroglia, neurolint, neurosci, nitrogen, notspecified, nursrep, nutraceuticals, nutrients, obesities, oceans, ohbm, onco, oncopathology, optics, oral, organics, organoids, osteology, oxygen, parasites, parasitologia, particles, pathogens, pathophysiology, pediatrrep, pets, pharmaceuticals, pharmaceutics, pharmacoepidemiology, pharmacy, philosophies, photochem, photonics, phycology, physchem, physics, physiologia, plants, plasma, platforms, pollutants, polymers, polysaccharides, populations, poultry, powders, preprints, proceedings, processes, prosthesis, proteomes, psf, psych, psychiatryint, psychoactives, psycholint, publications, purification, quantumrep, quaternary, qubs, radiation, reactions, realestate, receptors, recycling, regeneration, religions, remotesensing, reports, reprodmed, resources, rheumato, risks, robotics, rsee, ruminants, safety, sci, scipharm, sclerosis, seeds, sensors, separations, sexes, signals, sinusitis, siuj, skins, smartcities, sna, societies, socsci, software, soilsystems, solar, solids, spectroscj, sports, standards, stats, std, stresses, surfaces, surgeries, suschem, sustainability, symmetry, synbio, systems, tae, targets, taxonomy, technologies, telecom, test, textiles, thalassrep, therapeutics, thermo, timespace, tomography, tourismhosp, toxics, toxins, transplantology, transportation, traumacare, traumas, tropicalmed, universe, urbansci, uro, vaccines, vehicles, venereology, vetsci, vibration, virtualworlds, viruses, vision, waste, water, wem, wevj, wild, wind, women, world, youth, zoonoticdis

%---------
% article
%---------
% The default type of manuscript is "article", but can be replaced by: 
% abstract, addendum, article, benchmark, book, bookreview, briefcommunication, briefreport, casereport, changes, clinicopathologicalchallenge, comment, commentary, communication, conceptpaper, conferenceproceedings, correction, conferencereport, creative, datadescriptor, discussion, entry, expressionofconcern, extendedabstract, editorial, essay, erratum, fieldguide, hypothesis, interestingimages, letter, meetingreport, monograph, newbookreceived, obituary, opinion, proceedingpaper, projectreport, reply, retraction, review, perspective, protocol, shortnote, studyprotocol, supfile, systematicreview, technicalnote, viewpoint, guidelines, registeredreport, tutorial,  giantsinurology, urologyaroundtheworld
% supfile = supplementary materials

%----------
% submit
%----------
% The class option "submit" will be changed to "accept" by the Editorial Office when the paper is accepted. This will only make changes to the frontpage (e.g., the logo of the journal will get visible), the headings, and the copyright information. Also, line numbering will be removed. Journal info and pagination for accepted papers will also be assigned by the Editorial Office.

%------------------
% moreauthors
%------------------
% If there is only one author the class option oneauthor should be used. Otherwise use the class option moreauthors.

%---------
% pdftex
%---------
% The option pdftex is for use with pdfLaTeX. Remove "pdftex" for (1) compiling with LaTeX & dvi2pdf (if eps figures are used) or for (2) compiling with XeLaTeX.

%=================================================================
% MDPI internal commands - do not modify
\newcommand{\add}[1]{\textcolor{blue}{#1}}
\newcommand{\del}[1]{\textcolor{red}{\sout{#1}}}
\firstpage{1} 
\makeatletter 
\setcounter{page}{\@firstpage} 
\makeatother
\pubvolume{1}
\issuenum{1}
\articlenumber{0}
\pubyear{2026}
\copyrightyear{2026}
%\externaleditor{Firstname Lastname} % More than 1 editor, please add `` and '' before the last editor name
\datereceived{ } 
\daterevised{ } % Comment out if no revised date
\dateaccepted{ } 
\datepublished{ } 
%\datecorrected{} % For corrected papers: "Corrected: XXX" date in the original paper.
%\dateretracted{} % For retracted papers: "Retracted: XXX" date in the original paper.
\hreflink{https://doi.org/} % If needed use \linebreak
%\doinum{}
%\pdfoutput=1 % Uncommented for upload to arXiv.org
%\CorrStatement{yes}  % For updates
%\longauthorlist{yes} % For many authors that exceed the left citation part
%\IsAssociation{yes} % For association journals

%=================================================================
% Add packages and commands here. The following packages are loaded in our class file: fontenc, inputenc, calc, indentfirst, fancyhdr, graphicx, epstopdf, lastpage, ifthen, float, amsmath, amssymb, lineno, setspace, enumitem, mathpazo, booktabs, titlesec, etoolbox, tabto, xcolor, colortbl, soul, multirow, microtype, tikz, totcount, changepage, attrib, upgreek, array, tabularx, pbox, ragged2e, tocloft, marginnote, marginfix, enotez, amsthm, natbib, hyperref, cleveref, scrextend, url, geometry, newfloat, captiion, draftwatermark, seqsplit
% cleveref: load \crefname definitions after \begin{document}

%%% ===== Enable the red color for comments (Remove before submission) =====
\makeatletter\@ifpackageloaded{pgfplots}{}{\usepackage{pgfplots}}\makeatother
\pgfplotsset{compat=1.16}
\usepackage{xcolor}
\usepackage{comment}
% ---- NK change tracking ----
\usepackage{xcolor}
\usepackage[normalem]{ulem}   % \sout for strikethrough; normalem keeps \emph as italics
\newif\ifshowchanges
\showchangestrue              % set to \showchangesfalse for the clean final version

% Deleted text: red strikethrough (disappears in the final version)
\newcommand{\nkdel}[1]{\ifshowchanges{\color{red}\sout{#1}}\fi}
% Added text: blue with "NK:" label (becomes normal black text in the final version)
\newcommand{\nkadd}[1]{\ifshowchanges{\color{blue}\textbf{NK:} #1}\else #1\fi}
% ----------------------------

% Command for red comments
\newcommand{\commentRed}[1]{\textcolor{red}{#1}}

% Command for orange comments
\newcommand{\commentVP}[1]{\textcolor{orange}{\textbf{VP:} #1}}

% Command for blue comments
\newcommand{\commentNK}[1]{\textcolor{blue}{\textbf{NK:} #1}}

% Command for green comments
\newcommand{\commentTD}[1]{\textcolor{green}{\textbf{TD:} #1}}

% Command for green comments
\newcommand{\commentGN}[1]{\textcolor{olive}{\textbf{GC:} #1}}

%=================================================================
% Please use the following mathematics environments: Theorem, Lemma, Corollary, Proposition, Characterization, Property, Problem, Example, ExamplesandDefinitions, Hypothesis, Remark, Definition, Notation, Assumption
%% For proofs, please use the proof environment (the amsthm package is loaded by the MDPI class).

%=================================================================
% Full title of the paper (Capitalized)
%\Title{Benchmarking Local LLMs with LM Studio for Research Abstract Generation}
\Title{A Retrieval-Augmented Generation Framework for Knowledge Management Using Local Large Language Models}

% MDPI internal command: Title for citation in the left column
\TitleCitation{A Retrieval-Augmented Generation Framework for Knowledge Management Using Local Large Language Models}

% Author Orchid ID: enter ID or remove command
\newcommand{\orcidauthorA}{0000-0000-0000-000X} % Add \orcidA{} behind the author's name
%\newcommand{\orcidauthorB}{0000-0000-0000-000X} % Add \orcidB{} behind the author's name

% Authors, for the paper (add full first names)
%\Author{Osama Hamed $^{1,}$*\orcidA{}, Nadeem Zaidkilani $^{2}$, Firstname Lastname $^{3}$ and Ghaleb Natour $^{1,}$*}
\Author{Osama Hamed $^{1,2,}$*, Nadeem Zaidkilani $^{3}$, Tasneem Duridi $^{4}$, Vibhav Patil $^{5}$ and Ghaleb Natour $^{6}$}

%\longauthorlist{yes}

% MDPI internal command: Authors, for metadata in PDF
\AuthorNames{Osama Hamed, Firstname Lastname, Firstname Lastname and Firstname Lastname}

% MDPI internal command: Authors, for citation in the left column, only choose below one of them according to the journal style
% If this is a Chicago style journal 
% (arts, genealogy, histories, humanities, jintelligence, laws, literature, religions, risks, socsci): 
% Lastname, Firstname, Firstname Lastname, and Firstname Lastname.

% If this is a APA style journal 
% (admsci, behavsci, businesses, econometrics, economies, education, ejihpe, games, humans, ijfs, journalmedia, jrfm, languages, psycholint, publications, tourismhosp, youth): 
% Lastname, F., Lastname, F., \& Lastname, F.

% If this is a ACS style journal (Except for the above Chicago and APA journals, all others are in the ACS format): 
% Lastname, F.; Lastname, F.; Lastname, F.
\isAPAStyle{%
       \AuthorCitation{Lastname, F., Lastname, F., \& Lastname, F.}
         }{%
        \isChicagoStyle{%
        \AuthorCitation{Lastname, Firstname, Firstname Lastname, and Firstname Lastname.}
        }{
        \AuthorCitation{Hamed, O.; Zaidkilani, N.; Duridi, T.; Patil, V.; Natour, G.}
        }
}

% Affiliations / Addresses (Add [1] after \address if there is only one affiliation.)
\address{%
$^{1}$ \quad Institute of Technology and Engineering, Forschungszentrum Jülich, Jülich, Germany; o.hamed@fz-juelich.de\\
$^{2}$ \quad Computer Systems Engineering Department, Palestine Technical University - Kadoorie, Tulkarm, Palestine; osama.hamed@ptuk.edu.ps \\
$^{3}$ \quad Computer Science Apprenticeship Program, Faculty of Information Technology and Artificial Intelligence, An-Najah National University, Nablus, Palestine; nadimkilani@najah.edu\\
$^{4}$ \quad Department of Computer Science, Palestine Technical University - Kadoorie, Tulkarm, Palestine; tasneem.duridi@ptuk.edu.ps\\
$^{5}$ \quad Institute of Technology and Engineering, Forschungszentrum Jülich, Jülich, Germany; v.patil@fz-juelich.de\\
$^{6}$ \quad Institute of Technology and Engineering, Forschungszentrum Jülich, Jülich, Germany; g.natour@fz-juelich.de
}

% Contact information of the corresponding author
%\firstnote{Corresponding address: Computer Systems Engineering Department, Palestine Technical University - Kadoorie, Palestine.}
\corres{Correspondence: o.hamed@fz-juelich.de}

% Current address and/or shared authorship
%\firstnote{Current address: Affiliation.}  
% Current address should not be the same as any items in the Affiliation section.

%\secondnote{These authors contributed equally to this work.}
% The commands \thirdnote{} till \eighthnote{} are available for further notes.

%\simplesumm{} % Simple summary

%\conference{} % An extended version of a conference paper

% Abstract (Do not insert blank lines, i.e. \\)
\abstract{
Many research institutions possess growing collections of scientific papers and technical
documents that staff need to search and query. However, processing these materials
using cloud-hosted Large Language Models (LLMs) is often a concern due to privacy,
confidentiality, and cost constraints. While Retrieval-Augmented Generation (RAG) offers
a promising solution for grounding model responses in private documents and mitigating
hallucinations, most current RAG systems rely on cloud-based models and index full-text
segments; this increases computational costs and often leads to the retrieval of segments
with low relevance to the query. Furthermore, many internal documents lack concise
summaries that could guide the retrieval process. It remains unclear whether locally
deployed Small Language Models (SLMs) can generate summaries of sufficient quality to
support accurate retrieval within secure environments. To address this gap, we propose
a RAG-based knowledge base that employs an ``abstract-first retrieval'' methodology
and operates entirely on local infrastructure.
% NK change 1: grammar ("we apply ... are utilized"); reviewer comment 16
Three locally deployed LLMs (Mixtral-8$\times$7B-Instruct, Qwen3-VL-8B-Instruct and
Meta-Llama-3.1-8B-Instruct) were used to generate the abstracts through LM Studio's
REST API.
The quality of the generated abstracts is then assessed against the gold-standard
abstracts (a subset of the arXiv dataset) using BLEU, ROUGE, and BERTScore metrics.
% NK change 2: comma splice; reviewer comment 16
In a comparative study and an ablation study, Llama-3.1-8B outperformed the
other models.
Based on these results, we developed a RAG-based KB system that adopts an abstract-first
retrieval method. This method indexes paper abstracts and prioritizes semantic searches
within a separate vector database (VDB) of abstracts in secure settings.
% NK change 3: add evaluation results, remove design statement; reviewer comments 1, 14, 15
Finally, we evaluated the RAG system with the Retrieval Augmented Generation
Assessment (RAGAS) framework, using a local LLM (Qwen2.5-14B) as the judge, on 80 questions
with reference answers generated by an LLM from the abstracts. Since synthetic questions
can be flawed, each question was first validated against its source abstract and filtered further with a human annotator. Twenty-eight questions (35.0\%)
were invalid, mostly because they were too general to identify a single document. On the
52 valid questions, the system reached a faithfulness of 0.892, an answer correctness of
0.689, a context precision of 0.935 and a context recall of 0.753, and it ranked the correct
abstract first for 94.2\% of questions. 
These results showed us that abstract-first retrieval
with local models supports accurate retrieval, and that synthetic evaluation questions must
be validated before RAG performance can be assessed reliably.
% In particular, we utilized the web interface to search the KB for documents stored as embeddings of the Llama-3.1-8B-generated abstracts.
%\commentred{Continues as per the steps provided: 1, 2, 3.} This abstract should give a pertinent overview of the work. We strongly encourage authors to use the following style of structured abstracts, but without headings: (1) Background: place the question addressed in a broad context and highlight the purpose of the study; (2) Methods: describe briefly the main methods or treatments applied; (3) Results: summarize the article's main findings; (4) Conclusions: indicate the main conclusions or interpretations.
}

%\abstract{This study evaluates locally deployed Large Language Models (LLMs) for generating research abstracts using LM Studio, an offline framework for running open-source models. We address the need for privacy-preserving and cost-effective text generation in academic research environments. We developed a benchmarking pipeline to evaluate three instruction-tuned models—Qwen 3 8B, LLaMA-3.1 8B, and Mixtral 8×7B—using 200 Computer Science and Machine Learning papers from arXiv. Using LM Studio's OpenAI-compatible API, we generated abstracts locally and assessed their quality against gold-standard abstracts using BLEU, ROUGE, and BERTScore metrics. Results show that all models produce coherent and contextually relevant summaries, with Qwen 3 8B achieving the best balance between conciseness and semantic fidelity. Our experiments confirm that local LLMs can match the performance of cloud-based systems while maintaining data security and hardware control. These findings demonstrate that LM Studio provides a practical platform for reproducible text generation research and enables comparative studies on model efficiency and output quality in academic applications.}

% Keywords
\keyword{Large Language Models (LLMs); Local LLMs; Text Generation; Text Summarization; Evaluation Metrics; BLEU; ROUGE-L; BERTScore; Vector Database; RAG; RAGAS} 

% The fields PACS, MSC, and JEL may be left empty or commented out if not applicable
%\PACS{J0101}
%\MSC{}
%\JEL{}

%%%%%%%%%%%%%%%%%%%%%%%%%%%%%%%%%%%%%%%%%%
% Only for the journal Diversity
%\LSID{\url{http://}}

%%%%%%%%%%%%%%%%%%%%%%%%%%%%%%%%%%%%%%%%%%
% Only for the journal Applied Sciences
%\featuredapplication{Authors are encouraged to provide a concise description of the specific application or a potential application of the work. This section is not mandatory.}
%%%%%%%%%%%%%%%%%%%%%%%%%%%%%%%%%%%%%%%%%%

%%%%%%%%%%%%%%%%%%%%%%%%%%%%%%%%%%%%%%%%%%
% Only for the journal Data
%\dataset{DOI number or link to the deposited dataset if the dataset is published separately. If the dataset shall be published as a supplement to this paper, this field will be filled by the journal editors. In this case, please submit the dataset as a supplement.}
%\datasetlicense{License under which the dataset is made available (CC0, CC-BY, CC-BY-SA, CC-BY-NC, etc.)}

%%%%%%%%%%%%%%%%%%%%%%%%%%%%%%%%%%%%%%%%%%
% Only for the journal BioTech, Fishes, Neuroimaging and Toxins
%\keycontribution{The breakthroughs or highlights of the manuscript. Authors can write one or two sentences to describe the most important part of the paper.}

%%%%%%%%%%%%%%%%%%%%%%%%%%%%%%%%%%%%%%%%%%
% Only for the journal Encyclopedia
%\encyclopediadef{For entry manuscripts only: please provide a brief overview of the entry title instead of an abstract.}

%%%%%%%%%%%%%%%%%%%%%%%%%%%%%%%%%%%%%%%%%%
% Different journals have different requirements. Please check the specific journal guidelines in the "Instructions for Authors" on the journal's official website.
%\addhighlights{yes}
%\renewcommand{\addhighlights}{%
%
%\noindent The goal is to increase the discoverability and readability of the article via search engines and other scholars. Highlights should not be a copy of the abstract, but a simple text allowing the reader to quickly and simplified find out what the article is about and what can be cited from it. Each of these parts should be devoted up to 2~bullet points.\vspace{3pt}\\
%\textbf{What are the main findings?}
% \begin{itemize}[labelsep=2.5mm,topsep=-3pt]
% \item First bullet.
% \item Second bullet.
% \end{itemize}\vspace{3pt}
%\textbf{What is the implication of the main finding?}
% \begin{itemize}[labelsep=2.5mm,topsep=-3pt]
% \item First bullet.
% \item Second bullet.
% \end{itemize}
%}
% ---- NK change tracking (robust version) ----
\usepackage{xcolor}
\makeatletter
\@ifpackageloaded{ulem}{}{\usepackage[normalem]{ulem}}
\@ifpackageloaded{tikz}{}{\usepackage{tikz}}
\makeatother
\usetikzlibrary{positioning,arrows.meta}
\newif\ifshowchanges
\showchangestrue   % set to \showchangesfalse for the clean final version
\DeclareRobustCommand{\nkdel}[1]{\ifshowchanges{\color{red}\sout{#1}}\fi}
\DeclareRobustCommand{\nkadd}[1]{\ifshowchanges{\color{blue}\textbf{NK:} #1}\else #1\fi}
\newenvironment{nkaddblock}{\ifshowchanges\par\color{blue}\noindent\textbf{NK: new text begins.}\par\fi}{\ifshowchanges\par\noindent\textbf{NK: new text ends.}\par\fi}
\DeclareRobustCommand{\nkcolor}{\ifshowchanges\color{blue}\fi}
% ---------------------------------------------

%%%%%%%%%%%%%%%%%%%%%%%%%%%%%%%%%%%%%%%%%%
\begin{document}

%%%%%%%%%%%%%%%%%%%%%%%%%%%%%%%%%%%%%%%%%%
% \setcounter{section}{-1} %% Remove this when starting to work on the template.
% \section{Using this Template}
% Note that the order and names of article sections may differ from the requirements of the journal (e.g., the positioning of the Materials and Methods section). Please check the instructions on the authors' page of the journal to verify the correct order and names. 
%%For any questions, please contact the editorial office of the journal or support@mdpi.com. For LaTeX-related questions please contact latex@mdpi.com.
%%\endnote{This is an endnote.} % To use endnotes, please un-comment \printendnotes below (before References).


\section{Introduction}

%%The introduction should briefly place the study in a broad context \commentRed{Done: P1} and highlight why it is important \commentRed{Done: P2}. It should define the purpose of the work and its significance.  The current state of the research field should be reviewed carefully and key publications cited. Finally, briefly mention the main aim of the work and highlight the principal conclusions. As far as possible, please keep the introduction comprehensible to scientists outside your particular field of research.

Large language models (LLMs) have been described by \citet{brousseau2025llms} as ``the most promising technology to come along since, well, the printing press.'' However, they have limitations, such as knowledge cutoffs and outdated information, hallucinations, and lack of access to private institutional data \cite{bratanic2025essential}.
A promising approach to mitigate these limitations is Retrieval-Augmented Generation (RAG), which grounds the answers of an LLM in documents retrieved from an external collection.
RAG can be integrated with both closed-source and open-source LLMs.
%LLMs fall into two main categories: (i) closed-source LLMs, which are hosted on the cloud; and (ii) open-source LLMs, which are hosted on local machines.
Although cloud-based LLMs dominate today, they expose private documents to external providers. Recently, several tools have emerged to make local inference possible, such as LM Studio and Ollama.
%DEL%This paper investigates the use of lightweight and locally deployed LLMs to create a RAG-based knowledge base (KB) system from academic institutional documents and project archives.
%DEL%\commentRed{One more supporting sentence!} We split this further into two steps: abstract generation, and RAG-based KB system.

Managing long-term research projects in large enterprise environments often leads to stacks of information. A prime motivating example of this case can be seen at the Institute of Technology and Engineering (ITE), at Forschungszentrum Jülich.
There, a vast amount of historical project files has become practically inaccessible.
%DEL:%\commentRed{ We participate in, oversee, and manage long-term research projects. As a result, in some cases, we have project's folder containing few thousands of files in various formats. This equates to a large number of files in total (approx. 219k in PDF format only), but without any search criteria applied.}
Knowledge-base systems (KBSs) can provide effective solutions to make these documents searchable.
KBSs are typically data-driven.
However, the emergence of RAG techniques makes it feasible to build RAG-based KBS.
%Furthermore, most of the files lack a documentation summary (or abstract), hindering methods like abstract-first retrieval \cite{aytar2024retrieval}.
%\commentRed{DEL: Consequently, this makes the search for documents from past projects like searching a needle in a haystack.}
RAG-based KBS employs different retrieval strategies, such as abstract-first retrieval \cite{aytar2024retrieval} and small-to-big retrieval.
Abstract-first retrieval requires a summary or an abstract for every document; however, in an institutional repository where most documents lack one, this strategy cannot be applied.
%DEL:%To fix this issue, we are not interested in solutions that alleviate the problem, such as the ones based on Relational Database Management Systems (DBMS). We are rather interested in creating solution based on the emerging Artificial Intelligence (AI) techniques, in particular Retrieval-Augmented Generation (RAG) and  vector databases (VDB).
%DEL:%This paper simulates this scenario and proposes a RAG-based solution that addresses the different perspectives of this issue.
%REU:%We split this further into two steps: abstract generation, and RAG-based KB system.

Thus, one specific problem we address in this paper is the lack of concise document representations that can support semantic retrieval in large institutional repositories.
%It is difficult to apply the abstract-first retrieval approach when documents lack summaries. Keyword-based retrieval, on the other hand, remains dependent on the availability of appropriate terms and metadata.
RAG provides a mechanism for retrieving relevant documents through vector representations. However, the quality of the retrieval process depends on the quality of the representations stored in the vector database. Additionally, these representations must be created from confidential institutional documents, which cannot be sent to cloud-based LLM services.
This paper simulates this scenario and assumes the absence of these representative summaries. It proposes a RAG-based solution that considers the different perspectives of this issue: abstract generation, RAG system implementation, and RAG system evaluation.

%Due to their importance in the abstract-first retrieval method.
This paper begins by investigating the use of lightweight and locally deployed LLMs to automate the generation of abstracts from extensive institutional archives. This step plays a vital role in creating a vector database (VDB), and is considered a prerequisite for enabling similarity searches of institutional documentation in a secure setting.
%Our proposed solution is to build a knowledge base (KB) with a search criteria that enables text similarity, a feature of VDB.
%We thus want to employ a VDB to store the embeddings of the abstracts or the summary of each paper, and to provide a RAG-based search function that filters the documents based on their similarity with the typed text.
%In other words, as a first-stage, we want to enable finding the relevant document(s) using the high-level summary or the abstract.
%After having a look at the project folders, it can be clearly seen that a huge portion of the documents lacks to a summary or abstract.
This work addresses three research questions (RQs):
%%\commentRed{"DEL: How to generate an abstract/ or a summary for all papers that have no abstract, and without employing a human factor?". Shall we use classical machine learning (ML) algorithms, deep learning, or pre-trained LLMs?}
(i) How do quantized SLMs (8B parameters) and locally deployed LLMs compare in  generating semantically accurate abstracts when deployed on local hardware?; (ii) What effect do prompt structures and decoding parameters have on the accuracy and brevity of local LLM-generated abstracts?; and (iii)
%\nkdel{Does the abstract-first retrieval method using locally generated summaries improves the correctness of grounded-responses from LLMs.}
We evaluate the retrieval and answer accuracy of a retrieval-augmented
generation (RAG) pipeline built on abstract-first retrieval, and examine
how the quality of LLM-generated evaluation questions shapes this
assessment. This focus on LLM-based pipelines is motivated by prior work
showing their advantage over classical methods; for instance,
\citet{munzir2025large} reported that LLMs outperform classical machine
learning (ML) and deep learning (DL) approaches. Inspired by this finding, our proposed solution for the abstract generation relies solely on LLMs, in particular on a subcategory of LLMs that runs on local hardware, known as on-device LLMs \citep{qu2025mobileedgeintelligencelarge}.
Utilizing these LLMs has the advantage that no document leaves the local network.
%DEL:%\commentRed{More precisely, we want to adopt a few LLMs that belong to local or small LLMs family and evaluate them using industry standards for abstracts generation. Having that done, we use LangChain \cite{chase2022langchain} to implement a RAG-based KB showcase that adopts abstract-first retrieval method.}

To evaluate the LLMs used for abstract generation, we follow the principle stated by \citet{brousseau2025llms}, who said that ``when evaluating a model, you need two things: a dataset and a metric.'' Following this principle, we evaluated the local LLMs for this task using a carefully selected sub-dataset of the arXiv Dataset \citep{arxivdataset2024}. We apply standard evaluation metrics to assess the quality of the generated abstracts and identify the best performing local model. Completing this part allows us to continue with the RAG–based KB implementation, which stores the generated abstracts as embeddings in a vector database.
%%we create a RAG–based KB that allows users to find relevant documents based on textual similarity.
%This study thus demonstrates a practical application of the evaluation framework suggested by  \citet{brousseau2025llms}, linking model performance to measurable outcomes in a real-world institutional repository.

A further gap concerns how such systems are evaluated. Writing questions with known answers by hand is slow, so RAG systems are increasingly evaluated using questions and reference answers generated by an LLM from the indexed documents, and scored with frameworks such as RAGAS~\cite{es2024ragas}. This practice assumes that every generated question is valid. However, an LLM that sees one document at a time can write questions that are clear next to that document but ambiguous within a collection, such as ``What is the primary goal of the research presented in this paper?'', or can produce a wrong reference answer. Such items lower the measured performance for reasons unrelated to the system. For the third RQ, or (iii) above, we therefore validate every generated question against its source abstract before the evaluation and test two hypotheses: \textbf{H1}, generated questions that do not name their topic reduce retrieval of the correct document in a multi-document corpus; and \textbf{H2}, removing invalid generated questions changes the RAGAS scores of the RAG pipeline significantly.

Overall, this study demonstrates a practical RAG application of the evaluation framework presented in \cite{brousseau2025llms}.
This has the advantage of linking model performance to measurable outcomes on the selected arXiv subset. The institutional repository provides just the motivating use case.
%DEL%\commentRed{ The findings show that locally deployed LLMs can generate coherent and informative summaries that support semantic search within the proposed RAG-based knowledge base, which is also superior to traditional keyword-based approaches}.
As inspired by \citet{aytar2024retrieval}, the use of automatically generated abstracts makes abstract-first retrieval possible for project documents that have no abstracts. This study also demonstrates the potential of local LLMs for privacy-preserving knowledge management.
The main contributions and findings of this paper can be summarized as follows:

\begin{itemize}[labelsep=2.5mm,topsep=-3pt]

\item Develop an automated abstract generation pipeline that makes use of three locally deployed LLMs: Mixtral-8$\times$7B-Instruct (Mixtral-8$\times$7B), Qwen3-VL-8B-Instruct (Qwen3-VL-8B), and Meta-Llama-3.1-8B-Instruct (Llama-3.1-8B), none of which requires buying API keys. We established a comparative study to evaluate the three models under the same settings. We used three evaluation metrics: BLEU, ROUGE and BERTScore, and found that Llama-3.1-8B achieved the strongest performance among the tested models under the evaluated settings.

\item Conduct an ablation study to examine the effect of prompt design on Meta's Llama-3.1 model. We tested the model using 50 and 100 articles and report our results using the same evaluation metrics provided in the comparative study. Our results showed that the fully structured, instruction-based (zero-shot) prompt generates abstracts comparable to the gold standard in terms of both lexical coverage and semantic similarity.

\item %\nkdel{Use VDB and RAG techniques to create a KB system that brings external knowledge to LLMs. This approach enables semantic search and retrieves the most relevant documents without leaving institutional repositories.}
Build a local RAG-based KB that indexes one abstract per document in a VDB, and evaluate it with the RAGAS framework using a local judge model (Qwen2.5-14B, with nomic-embed-text embeddings and top-20 retrieval). On 52 validated questions, the system ranked the correct abstract first for 94.2\% of questions, retrieved it in the top 20 for 98.1\%, and reached a faithfulness of 0.892, an answer correctness of 0.689, a context precision of 0.935 and a context recall of 0.753. 
Context recall proved unreliable with the Qwen2.5-14B as judge model. Thus, retrieval was also measured directly from document identifiers (i.e., Hit@1 and Hit@20).

\item Furthermore, create a labeled dataset for abstract generation consisting of 200 papers from the papers
submitted to arXiv in April 2007 (identifiers 0704.0002 to 0704.1627) and covering topics from  physics, mathematics, biology and computer science. 
This sub-dataset was extracted from the arXiv Dataset and was used as the ground truth to evaluate the three local LLMs in this study. The abstracts generated for these 200 papers form the corpus of the RAG-based KB.

\item Show that LLM-generated evaluation questions must be validated before use: 28 of 80 generated questions (35.0\%) were invalid, mostly because they were too general to identify their source document (25 of 28). The deleted questions scored significantly lower than the kept ones on three of the four RAGAS metrics (two-sided Mann--Whitney U test, $p \le 0.001$); removing them raised answer correctness from 0.569 to 0.689, which supports H2. General questions retrieved their source document in only 40.0\% of cases, against 98.1\% for valid questions, which supports H1.

\end{itemize}\vspace{3pt}
%%%%%%%%%%%%%%%%%%% DONT DELETE %%%%%%%%%%%%%%%%%%%%%%%
%Nevertheless, some limitations remain. The evaluation was performed on documents from a single domain, context, limiting generalizability. The computational cost of running local LLMs may affect scalability for very large datasets. Additionally, while the generated summaries are generally accurate, some domain-specific details may be missing, suggesting a need for further refinement or human-in-the-loop validation.

%The remainder of this paper is structured as follows: Section~\ref{sec:related_work} reviews related studies; Section~\ref{sec:methodology} outlines the dataset, methodology, evaluation metrics, and experimental setup; Section~\ref{sec:results} presents and discusses the results; and finally Section~\ref{sec:conclusion} draws conclusions and suggestions for future research directions.

%%%%%%%%%%%%%%%%%%%%%%%%%%%%%%%%%%%%%%%%%%
% \section{Background}
% \textbf{NZ: Add your text here!}
\begin{comment}
\section{Related Work}
\label{sec:related_work}

Automatic text summarization has a long history, with early work dating to \citep{luhn1958automatic} that introduced term frequency-based statistical methods for sentence selection. Today's text summarization has fully adopted neural and transformer-based methods for generation. The introduction of large language models (LLMs) like GPT-3 and LLaMA have greatly improved summarization. Unlike previous models that generated either extractive or abstractive summaries, LLMs can do both and produce accurate, integrated summaries from little input. BART and T5 are examples of encoder-decoder models which work well for summarization because their architecture allows for effective use of structured external knowledge sources \citep{daraghmi2025comparative}.

The connection between the quality of the summarization and the performance of the retrieval system is already established. \citep{bramer2018searching} showed the impact of keyword coverage in the document abstracts and its effect on search recall in MEDLINE, stating that the quality of the summary is not the most important aspect of its summary, because it has an effect on the retrieval of the documents. \citep{ketcham2010editors} pointed out that good abstracts are the best entry points for indexing and retrieval.These findings are directly relevant to our work, where generated abstracts serve not only as summaries but as the semantic representations that drive retrieval within a vector database.

The use of vector databases (VDBs) for semantic document retrieval has gained significant traction in recent years. Retrieval-Augmented Generation (RAG), introduced by Lewis et al.~\cite{lewis2020retrieval}, combines pre-trained parametric memory (i.e., large language models) with non-parametric external knowledge stored in a dense vector index, enabling dynamic access to domain-specific information at query time. As a result, RAG has emerged as a predominant architecture for grounding language model outputs in organizational and domain-specific document collections~\cite{lewis2020retrieval}. Within this framework, documents, or their corresponding summaries, are encoded as dense vector embeddings and stored in a VDB, which is queried using similarity measures such as cosine similarity or dot-product similarity to retrieve the most relevant contextual information for a given input. Consequently, the quality of generated abstracts becomes a critical factor, as semantically weak or imprecise summaries yield poor embeddings, thereby degrading downstream retrieval performance across the entire knowledge base. Motivated by this dependency, this work approaches abstract generation not merely as a standalone text generation task, but as a foundational component that directly influences the effectiveness of document retrieval in RAG-enabled institutional knowledge systems.

An early and influential study on LLM-based abstract generation is \citep{gao2023comparing}, who evaluated ChatGPT’s ability to produce scientific abstracts for medical journal papers, comparing them against human-written originals using AI detectors and blinded expert reviewers. Their finding that generated abstracts were largely indistinguishable from real ones — with human reviewers correctly identifying only 68\% of generated abstracts — established the viability of LLM-based abstract generation as a serious research direction. However, their study used a cloud-based LLM (GPT-3), generated abstracts for existing papers that already had reference summaries, and was not concerned with downstream retrieval.  

 A recent study \citep{bektacs2025academic} conducted a benchmark analysis to evaluate how effectively different large language models (LLMs) can generate abstracts of academic articles. The authors examined nine models, including large proprietary systems such as Gemini-2.0-Flash and Mistral-Large-2411, as well as smaller open-source models such as LLaMA-3.1-8B, DeepSeek-R1-8B, Qwen2.5 and Gemma-2. Using a dataset of 1,500 English papers collected from the ACL Anthology, the models were prompted in a zero-shot setting to generate abstracts as expert scientific writers. The generated outputs were evaluated using several automatic metrics, including ROUGE, BERTScore, cosine-based sentence embedding similarity, UniEval, and G-Eval. The results showed that larger models generally produced abstracts that were semantically closer to the original ones, while semantic metrics provided more reliable comparisons than simple word-overlap measures. However, the study relied solely on prompt engineering and publicly available data, without grounding generation in domain-specific or private document collections, a limitation addressed in our RAG-based approach.

The most closely related previous study is \citep{kim2025exploring}, which systematically evaluated LLM-generated abstracts for academic papers from the ASIS\&T 2024 conference, measuring semantic similarity using BERT-based cosine similarity and collecting expert judgments. Their finding that domain experts preferred AI-generated abstracts when semantic discrepancies existed with author-written ones provides important empirical grounding for using LLM-generated summaries in real-world settings. However, their study focused on cloud-based LLMs and academic publications. Our work differs by applying local, quantized open-source LM Studio to institutional project documents that lack summaries, and by directly evaluating the retrieval effectiveness of generated abstracts within an RAG-enabled knowledge base.


\begin{table}[H]
\centering
\caption{Comparison with Recent Generation and Deployment Studies}
\label{tab:comparison_recent}
\renewcommand{\arraystretch}{1.2}
\setlength{\tabcolsep}{4pt}
\footnotesize

\begin{tabular}{p{2cm} p{3.6cm} p{2.6cm} p{4 cm} p{1.2cm}}
\toprule
\textbf{Study} & \textbf{Task / Setting} & \textbf{Model} & \textbf{RAG + Evaluation} & \textbf{Privacy} \\
\midrule

\citep{kim2025exploring} 
& Academic abstract generation 
& GPT 
& No RAG; Cosine similarity, expert evaluation 
& No \\

\citep{gao2023comparing}
& Medical abstract generation 
& GPT-3 (ChatGPT) 
& No RAG; AI detector and human evaluation 
& No \\

\citep{bektacs2025academic} 
& Academic abstract generation with multi-model benchmarking 
& Gemini, Mistral, LLaMA, DeepSeek, Qwen 
& No RAG; ROUGE, BERTScore, UniEval, G-Eval 
& No \\

This Study 
& Abstract generation with RAG over local documents 
& LLaMA, Mistral
&RAG; BLEU, ROUGE, BERTScore 
& Yes \\

\bottomrule
\end{tabular}

\end{table}
\end{comment}
\section{Related Work}
\label{sec:related_work}

\subsection{Why Missing Abstracts Matter for Retrieval?}

A RAG system can only answer well if its retriever first finds the right documents~\citep{lewis2020retrieval}. In dense retrieval, queries and documents are mapped to vectors, and relevance is judged by their similarity~\citep{karpukhin2020dense,reimers2019sentence}. 
What the system retrieves therefore depends heavily on what text was embedded for each document in the first place.

For scholarly material, the title and abstract have long served this purpose. An abstract states the problem, method and main result in a few hundred words, which is close to ideal for a single document vector. SPECTER builds its scientific document embeddings from exactly these two fields~\citep{cohan2020specter}. Without an abstract, there is no obvious replacement. The full text often exceeds the input length of common embedding models, so it has to be split into passages, and the document is then represented by many local vectors rather than one global one. Language models also make poor use of information buried in the middle of long inputs~\citep{liu2024lost}, and dense retrievers are known to transfer unevenly across domains and document types~\citep{thakur2021beir}.

This is not a marginal case. Institutional repositories often contain technical reports, theses, project deliverables and internal documents that were never written with an abstract. For these collections, the question is not which retriever to use, but what text should represent each document at all.

\subsection{Generated Text as a Stand-In Representation}

Several retrieval methods already rely on text produced by a model rather than by an author. Doc2query predicts the queries that a document might answer and adds them to the index~\citep{nogueira2019document}. HyDE works from the other side, writing a hypothetical answer document for each query and searching with its embedding~\citep{gao2023precise}. RAPTOR summarizes chunks recursively and arranges the summaries in a tree for long-document question answering~\citep{sarthi2024raptor}. Each of these shows that generated text can fill gaps in a weak representation. None of them, however, asks whether a generated abstract can stand in for a missing author abstract at the level of the whole document.

That use carries its own risk. Abstractive summarizers often produce statements that the source does not support~\citep{maynez2020faithfulness}. In a knowledge base, such an abstract does not merely mislead a reader; it becomes the document's identity in the vector space, and every later search is shaped by it. For this reason, faithfulness to the source matters at least as much as fluency.

\subsection{LLMs and Scientific Abstract Generation}

\citet{gao2023comparing} show how far fluency and faithfulness can diverge. They gave ChatGPT only the title and target journal of 50 papers from five leading medical journals and asked it to write the abstracts. Blinded reviewers recognized 68\% of these as machine-written, yet also labeled 14\% of the genuine abstracts as generated, and An AI-output detector (the GPT-2 Output Detector) separated generated from genuine abstracts with an area under the receiver operating characteristic curve (AUROC) of 0.94. Because the model never saw the papers, the abstracts reported invented results, and only 16\% followed the journal's required format. The study shows that an abstract can be convincing without being true to its paper. For retrieval, that is the failure that matters most.

\citet{bektacs2025academic} addressed this by generating abstracts from the papers’ own section text. Working with 1,500 ACL Anthology papers, they compared nine models, ranging from 1.5B-parameter open models to Mistral-Large-2411 (123B) and Gemini-2.0-Flash, using one structured zero-shot prompt. Evaluation combined embedding similarity, BERTScore, ROUGE, UniEval and G-Eval. The two large models produced the abstracts closest to the originals, models between 2B and 8B parameters (including Llama-3.1-8B) performed at a similar intermediate level, and the two smallest models performed worst. BERTScore separated the models only slightly, whereas ROUGE showed clearer differences. UniEval, a learned evaluator that scores summaries on coherence, consistency, fluency and relevance \citet{zhong2022towards}, rated consistency with the source, that is, factual alignment, as the weakest dimension for every model. The study did not report how the models were deployed, did not examine quantization, tested a single prompt, and did not use the abstracts for retrieval.

Prompt design is a further open point. \citet{kim2025exploring} generated abstracts with GPT-4o for 120 papers from the ASIS\&T~2024 conference, supplying the full text with the authors' abstracts and references removed, and compared six prompting strategies. The prompts changed the length and wording of the abstracts, but BERTScore against the authors' abstracts stayed almost constant (0.87--0.88), indicating no meaningful semantic difference. Notably, most generated abstracts were lexically closer to the full text than the authors' own abstracts were, and for the two papers whose author abstracts diverged most from the generated ones, seven of eight experts preferred the generated version. 
\citet{zhang-etal-2024-benchmarking} likewise found that zero-shot news' summaries from large models were judged by comparison with those written by freelance writers.
This suggests that frontier models need less prompt engineering for summarization. Whereas, smaller models follow instructions less reliably. 
Whether the same insensitivity holds for them still need to be tested.

\subsection{Local Deployment and Quantization}

Many organizations cannot send internal documents to an external API. For them, the realistic option is an open-weight model running on their own hardware, usually quantized to fit in limited memory. Post-training quantization methods such as GPTQ reduce memory and inference cost while preserving much of a model's accuracy~\citep{frantar2022gptq}, but the loss is not uniform: it grows at lower bit widths and varies across tasks and model families~\citep{jin2024comprehensive}. Combined with the gap between large and small models reported by~\citet{bektacs2025academic}, this leaves a practical question unanswered. Can small dense models of about 8B parameters, or sparse mixture-of-experts models that still fit on a single workstation, produce abstracts good enough to represent a document once quantized?

Local RAG systems have been built before. \citet{krutli2025augmenting} combined Mixtral-8$\times$7B with a Milvus vector database to answer questions from corporate real-estate and facility-management documents, and evaluated the system on 30 question--answer pairs using ROUGE, BLEU and semantic similarity. That evaluation targeted the generated answers; the quality of the document representations behind them was not examined separately.

\subsection{Evaluating Generated Abstracts}

BLEU and ROUGE count shared words and phrases, so they reward close wording and penalize valid paraphrase. Large-scale human evaluation has shown that they agree only weakly with expert judgments of summary quality~\citep{fabbri2021summeval}. BERTScore instead compares contextual token embeddings and captures meaning beyond exact matches~\citep{zhang2019bertscore}. \citet{bektacs2025academic} reached the same conclusion about the value of semantic metrics. Since dense retrieval also works on embedding similarity, semantic closeness to the author's abstract is a more relevant signal than lexical overlap of whether a generated abstract can serve as a retrieval representation. We therefore report BERTScore together with BLEU, ROUGE-1 and ROUGE-L (Section~\ref{sec:metrics}).

% NK change 4: corpus is not limited to CS/ML (arXiv IDs 0704.0002-0704.1627); reviewer comment 4
Overall, this study addresses the research questions (RQs) mentioned in the introduction. It uses 200 arXiv papers from April 2007 covering several fields (Subsection~\ref{sec:dataset}). 
Table~\ref{tab:comparison_recent} compares this work with the closest prior studies.

\begin{table}[H]
\footnotesize
\centering
\caption{Comparison of the present study with related work on LLM-based abstract generation and local RAG systems.}
\label{tab:comparison_recent}
\renewcommand{\arraystretch}{1.25}
\setlength{\tabcolsep}{3pt}
\footnotesize
\begin{tabular}{p{1.8cm} p{2.5cm} p{2.4cm} p{2.7cm} p{1.35cm} p{0.9cm}}
\toprule
\textbf{Study} & \textbf{Data and input} & \textbf{Model(s)} & \textbf{Evaluation} & \textbf{Retrieval use} & \textbf{Local} \\
\midrule
\citet{gao2023comparing} & 50 medical papers; title and journal only & ChatGPT & AI detection, plagiarism checks, blinded reviewers & No & No \\[3pt]
\midrule
\citet{bektacs2025academic} & 1{,}500 ACL papers; section text & 9 LLMs, 1.5B--123B & Embedding similarity, BERTScore, ROUGE, UniEval, G-Eval & No & Not reported \\[3pt]
\midrule
\citet{kim2025exploring} & 120 ASIS\&T papers; full text; six prompting strategies & GPT-4o & BERTScore, TF-IDF cosine, expert review & No & No \\[3pt]
\midrule
\citet{nogueira2019document,gao2023precise} & Retrieval collections; generated queries or hypothetical documents & Seq2seq models / LLMs & Retrieval metrics & Query-level text & -- \\[3pt]
\midrule
\citet{krutli2025augmenting} & Corporate real-estate and facility-management documents & Mixtral-8$\times$7B & ROUGE, BLEU, semantic similarity on 30 Q\&A pairs & No & Yes \\[3pt]
\midrule
% NK change 5: "This study" row updated to match the corpus and the RAG evaluation (reviewer comments 2, 4)
\textbf{This study} & \textbf{200 arXiv multi-field papers; full text} & \textbf{Llama-3.1-8B, Qwen3-VL-8B, Mixtral-8$\times$7B; quantized} & \textbf{BLEU, ROUGE-1/L, BERTScore; prompt ablation plus RAGAS (4 metrics) and Hit@$k$ on validated synthetic questions} & \textbf{Yes; local VDB} & \textbf{Yes} \\
\bottomrule
\end{tabular}
\end{table}

%%%%%%%%%%%%%%%%%%%%%%%%%%%%%%%%%%%%%%%%%%
\section{Materials and Methods}
\label{sec:methodology}
This section briefly describes the dataset used, LM Studio, local LLMs, evaluation metrics, and the experimental setup.

%Materials and Methods should be described with sufficient details to allow others to replicate and build on published results.

% \subsection{GenAI Usage} 
% In this section, where applicable, authors are required to disclose details of how gen-erative artificial intelligence (GenAI) has been used in this paper (e.g., to generate text, data, or graphics, or to assist in study design, data collection, analysis, or interpretation).

\subsection{Dataset}
\label{sec:dataset}

This study utilizes the arXiv Dataset \citep{arxivdataset2024}, publicly available on Kaggle\footnote{https://www.kaggle.com/datasets/Cornell-University/arxiv} platform. The arXiv Dataset is a repository of more than 1.7 million articles. The dataset, which is updated periodically, consists primarily of a large JSON metadata file (approximately 5 GB) containing key paper details with relevant features such as article titles, authors, categories, abstracts, full-text PDFs, and more. 

\paragraph{Data Preprocessing and Preparation}
%\subsubsection{Data Preprocessing and Preparation}
For the experiments conducted in this paper, we are not going to use all of the data. Instead, we are only going to experiment with a subset of arXiv Dataset, which suits the locally acquired computational power. We have implemented a data pipeline that reads the JSON line by line, and filters the papers based on a specific criterion.
% NK note for VP/OH: reviewer comment 4 asks for the exact filtering criteria (arXiv categories, time range, minimum text length, duplicate handling) and why 200 papers are representative; please state them here.
 The whole process is depicted in Figure~\ref{fig:Prep-prep}. We picked up a reasonable subset of 200 papers that include the full-paper text 
% NK change 1: the 200 papers (arXiv IDs 0704.0002-0704.1627) cover several fields, not only CS/ML; reviewer comment 4. Team: please confirm the field list.
and were submitted to arXiv in April 2007 (identifiers 0704.0002 to 0704.1627). These papers cover topics from physics, mathematics, biology and computer science. Figure~\ref{fig:articles} demonstrates how the 200 articles were selected. We stored the matching papers in a CSV file with four columns corresponding to the following header: \textbf{[\text{paper\_id}, title, abstract, body]}.

\begin{figure*}[ht]
\centering
\includegraphics[width=0.95\textwidth, height=0.53\textheight]{img/filter.png}
% NK change 2: caption grammar; reviewer comment 16
\caption{Selection of the 200 articles.}
\label{fig:articles}
\end{figure*}

%TODO: OH, switch figures.
\begin{figure*}[ht]
\centering
\includegraphics[width=0.95\textwidth, height=0.32\textheight]{img/napkin.png}
\caption{Data preprocessing workflow (from JSON to CSV).}
\label{fig:Prep-prep}
\end{figure*}

\subsection{Methodology}
% NK change 3: third stage added for the evaluation (reviewer comments 3, 5)
In order to develop the RAG-based KB system, we divided our work into three main stages: (i) abstract generation that carried out on LM Studio; (ii) a RAG-based KB system; and (iii) the evaluation of the KB system with validated synthetic questions (Subsection~\ref{sec:rag_eval}).  
First, we generated the abstracts using small, locally deployed, and pretrained LLMs. The LLMs are utilized as embedded models within the LM Studio~\footnote{The LM Studio command is typically used to start a local inference server for loaded models, e.g., \textbf{lmstudio serve --port 1234}}, an open-source platform for running LLMs locally. Then, as part of the RAG pipeline, % NK change 4: grammar ("stored to"); reviewer comment 16
we stored the generated abstracts in the VDB. Finally, end users are capable of retrieving the matching documents from the KB system based on a similarity search. Figure~\ref{fig:system_architecture} depicts the KB system architecture. 
%%Next, we present the LM Studio and the used local LLMs.

\begin{figure*}[ht]
\centering
\includegraphics[width=0.95\textwidth]{img/system_architecture.png}
\caption{Proposed KB system architecture divided into three main stages: processing, generation, and retrieval.}
\label{fig:system_architecture}
\end{figure*}

%Moved: After §Local LLMs % \subsection{Evaluation of the RAG-based KB}
% NK change 5: new subsection with the RAG evaluation method, Figure and Table (reviewer comments 3, 5, 6, 9, 11)

\subsubsection{LM Studio}
%%\commentRed{Relocate} 
LM Studio is a robust cross-platform desktop application compatible with macOS, Linux, and Windows. It's engineered to enable local execution of LLMs, providing developers and researchers with direct access to advanced model capabilities without depending on remote servers.

At its core, LM Studio functions as an integrated environment for downloading, managing, and running open-source LLMs such as Llama 3.1, Phi-3, and Gemma 2. It supports model quantization, custom inference configurations, and performance optimization for different hardware setups.

%% TODO: OH
%% Em dash: Parenthetical break!
LM Studio bridges the gap between local and cloud-based AI workflows. It allows users to experiment with, benchmark, and deploy models efficiently while maintaining full control over data, performance, and privacy.

Our abstract generation experiments were carried out using LM Studio \textit{Version 0.3.31} or higher, which was downloaded from the official LM Studio website~\footnote{https://lmstudio.ai/download} and installed on a Windows machine.

\subsubsection{Local LLMs}
The local LLMs play a vital role in our secure settings for abstracts generation and semantic search.
We implemented several NLP pipelines that call the LM Studio's REST compatible APIs, which are also compatible with the OpenAI API format. Our implementation utilizes the following families of local LLMs to automatically generate abstracts. All of these local LLMs are packaged using the GGUF format. GGUF\footnote{It replaced the older GGML (Georgi Gerganov Machine Learning) format in 2023 and has become the primary format for distributing quantized LLMs.} stands for "GGML Unified Format".
These models are summarized in Table~\ref{tab:local_LLMs} for quick reference. More details, such as the exact model quantizations, are provided in Appendix~\ref{app:sub_sw_hw}.

\begin{itemize}
\item	Mixtral-8$\times$7B-Instruct (Mixtral-8$\times$7B): % NK fix (factual): Mixtral-8$\times$7B is a sparse mixture-of-experts model with 46.7B parameters in total, of which 12.9B are used per token (Mistral AI model release); it is not a 7B model. Reviewer comment 7.
This is a sparse mixture-of-experts model of the Mixtral 8$\times$7B family, with 46.7B parameters in total, of which 12.9B are used per token. This model is treated as a baseline. Mixtral-8$\times$7B represents the technical name.
\item	Qwen3-VL-8B-Instruct (Qwen3-VL-8B): This is an 8 billion (8B) parameters model, which belongs to Qwen3-VL family, from Alibaba’s AI. Note that Qwen3-VL-8B and Qwen3-VL:8B refer to the same model. % NK fix (grammar): plural subject
Qwen3-VL:8B refers to the format used by local deployment tools like LM Studio and Ollama. 
% NK change 1: name the model used in the RAG evaluation; reviewer comments 3, 6
In the RAG evaluation stage, Qwen2.5-14B (qwen2.5:14b), run through Ollama, serves as the judge model of the RAGAS framework (Subsection~\ref{sec:rag_eval}).
\item	Meta-Llama-3.1-8B-Instruct (Llama-3.1-8B): This is an 8B parameters model that belongs to LLaMA-3.1 family. Llama-3.1-8B is the official technical name, whereas Llama3.1:8B  format is typically used by LM Studio and Ollama. % NK fixes (grammar): "This An" -> "This is an"; "where as" -> "whereas"
% NK change 2: name the model used in the RAG evaluation; reviewer comments 3, 6, 11
In the RAG evaluation stage, Llama 3.1 (llama3.1), run through Ollama, generates the evaluation questions and their reference answers (Subsection~\ref{sec:rag_eval}).
\end{itemize}

\begin{table}[H]
\centering
\caption{Overview of the local large language models (LLMs) used in this study.}
\label{tab:local_LLMs}
\renewcommand{\arraystretch}{1.2}
\begin{tabularx}{\textwidth}{llcX}
\toprule
\textbf{Model} & \textbf{Format} & \textbf{Model Size} & \textbf{Training Data} \\
\midrule
Mixtral-8$\times$7B-Instruct & GGUF*
& 46.7B 
& Proprietary dataset extracted from the open web. \\
\midrule
Qwen3-VL-8B-Instruct & GGUF
& 8B 
& Massive, highly diverse, mixture of datasets tailored for both visual and linguistic understanding. \\
\midrule
Meta-Llama-3.1-8B-Instruct & GGUF
& 8B 
& About 15 trillion tokens from publicly available sources. \\
\bottomrule
\end{tabularx}
\noindent{\footnotesize{\textsuperscript{1} It stands for GGML Unified Format.}}
\end{table}

\subsubsection{Experimental Pipelines}
Our pipelines (abstracts generation and RAG evaluation) are written in Python and are intended to run in a Python development environment with Python \textit{version 3.11.9}, or higher. % NK change 3: execution environment of the RAG evaluation; reviewer comments 3, 12

%\paragraph{Abstracts Generation}
On the one hand, all of the abstract generation pipelines use zero-shot prompting, which proves that text summarization can be effectively handled via simple prompts and no training examples \cite{zhang-etal-2024-benchmarking}. However, the first NLP pipeline, i.e. the \textbf{Mixtral-8$\times$7B-Instruct}'s model pipeline utilizes a long-prompt because Mixtral model was irresponsive with short prompts. The remaining pipelines are utilizing the same short-prompt, see Table~\ref{tab:prompt}, which instructs the local LLMs to generate an abstract of a specific word length. In our case, we opted for an abstract of 120--130 words, as inspired by the calculated average word length of 125 words. % NK fixes (grammar): "instruct" -> "instructs"; missing articles; "word" -> "words"; "Here in and after" -> "Hereafter"; "include" -> "includes"
Hereafter, the \textit{average word length} is denoted as (\textbf{Avg. W. Len}). All the generated abstracts are persisted to CSV files (having the header: \textbf{ [\text{paper\_id},  \text{gen\_abstract}]}) for later reference. This includes, among others, evaluating the LLMs using the adopted evaluation metrics (see Subsection~\ref{subsec:eval}).
%\paragraph{RAG Evaluation}

On the other hand, the abstract generation experiments use locally served language models rather than API keys for external model. However, the RAG evaluation pipeline, runs in Google Colab on an NVIDIA Tesla T4 GPU (Table~\ref{tab:rag_eval_setup}). Therefore, this evaluation environment is not a fully institution-hosted deployment. While the evaluation models were served locally through Ollama within the execution environment, this setup alone does not ensure zero data egress from the institution.
Hence, this configuration demonstrates the feasibility of using locally deployed models for RAG evaluation. Whereas, the deployment under strict institutional data-governance requirements still requires further verification.

% NK note (no change made, style only): "which proves that" may overstate what one citation shows (e.g., "since prior work shows that"); "irresponsive" is rare, "unresponsive" is more usual.

\begin{table}[H]
\caption{Prompts used for abstract generation across all evaluated local LLMs.\label{tab:prompt}}
\begin{tabularx}{\textwidth}{lX}
\toprule
\textbf{Prompt Type}& \textbf{Prompt Text} \\
\midrule
Long & Write ONLY a new formal academic abstract of exactly 120--130 words in one continuous paragraph.
Strictly cover: (1) background or motivation, (2) methods or approach,
(3) main results, (4) conclusions or implications.
Do NOT include feedback, suggestions, critiques, numbered lists, or references to ``your abstract''.
Base the abstract solely on the provided paper content. \\
\midrule
Short & Write a formal academic abstract of about 120--130 words as one paragraph.
Cover background, methods, main results, and conclusions clearly. \\%% Like P4
\bottomrule
\end{tabularx}
% \noindent{\footnotesize{\textsuperscript{1} Short prompts keep the benchmark fair and reproducible across models.}}
\end{table}

% \paragraph{- Qwen3-VL-8B-Instruct  Qwen 3 8B} qwen3-vl-8b-instruct
% \paragraph{- LLaMA-3.1 8B}
% \paragraph{- Mixtral 8×7B}

%% ------ ------ Begin of § , was 3.3 ------ ------
\subsection{Evaluation of the RAG-based KB}
\label{sec:rag_eval}
Figure~\ref{fig:rag_eval} summarizes how the KB system was evaluated, and Table~\ref{tab:rag_eval_setup} lists the settings.


\begin{figure*}[ht]
\centering
\includegraphics[width=0.95\textwidth]{img/rag_evaluation_workflow.pdf}
\caption{The evaluation workflow of the RAG-based KB. Questions and reference answers are generated from a random sample of the indexed abstracts and validated against their source abstracts; RAGAS and Hit@$k$ scores are reported for the valid questions.}
\label{fig:rag_eval}
\end{figure*}


\paragraph{Indexing and retrieval}
The KB indexes the 200 generated abstracts, one per paper. Each abstract is stored as a single chunk: abstracts are not split, so there is no chunk overlap, and chunk lengths range from 106 to 233 words (mean 137.6). Before indexing, the introductory sentences that the generation model added to its output (for example, ``Here is a combined abstract\ldots'') were removed; this was the only preprocessing step. For each question, the 20 most similar abstracts (top-$k = 20$) are retrieved and passed, together with their paper identifiers, to the answer-generation LLM. 
The answer and the 20 retrieved contexts were stored to serve as the basis for the evaluation.

\paragraph{Statistical analysis}
For each metric, we report the mean, standard deviation and 95\% confidence interval of the mean obtained by bootstrap resampling of questions (10{,}000 resamples).
%DEL: Not here%\commentVP{and also perform paired t-tests to check the statistical significance}. 
Kept and deleted questions were compared with the two-sided Mann--Whitney U test, and retrieval rates with Fisher's exact test ($\alpha = 0.05$). The evaluation was run once and the judge model used a temperature of 0.

\begin{table}[H]
\centering
\caption{The evaluation settings for RAG.}
\label{tab:rag_eval_setup}
\renewcommand{\arraystretch}{1.2}
\footnotesize
\begin{tabular}{p{0.25\linewidth} p{0.68\linewidth}}
\toprule
\textbf{Component} & \textbf{Setting} \\
\midrule
Indexed corpus & 200 generated abstracts; one abstract per chunk; no splitting, no overlap. \\
\midrule
Retrieval & top-$k = 20$ abstracts per question, passed with their paper identifiers to the answer-generation LLM. \\
\midrule
Question generation & Llama~3.1 (llama3.1) through Ollama; 80 abstracts sampled at random; one question and reference answer each. \\
\midrule
Evaluation framework & RAGAS 0.4.3: Faithfulness, AnswerCorrectness, ContextPrecision, ContextRecall; default metric settings. \\
\midrule
RAGAS judge model & Qwen2.5-14B (qwen2.5:14b) through Ollama's OpenAI-compatible API (openai 1.109.1); temperature 0; maximum 4{,}096 output tokens. \\
\midrule
RAGAS embedding model & nomic-embed-text through Ollama; used by answer correctness. \\
\midrule
Hardware & Google Colab, NVIDIA Tesla T4 GPU (15{,}360 MiB); all models run locally. \\
\bottomrule
\end{tabular}
\end{table}
%% ------ ------ End of § ------ ------

\subsection{Evaluation Metrics}
\label{sec:metrics}

To assess the quality of the generated outputs, we must select the appropriate evaluation metrics. To do so, we need to determine the broad category of abstract generation in the field of NLP. % NK fixes (typos): "brand" -> "broad"; "filed" -> "field"
 In NLP, automatic abstract generation is typically modeled as a specialized text summarization task, and text summarization itself is widely categorized under natural language generation (text generation) within NLG/NLG-style sequence‑to‑sequence models \citep{becker2024textgenerationsystematicliterature}. Thus, we employ several widely used automatic evaluation metrics in NLG, including BLEU, ROUGE, and BERTScore. These metrics provide complementary perspectives by evaluating both lexical overlap and semantic similarity between the generated text and the reference texts.

\textbf{BLEU (Bilingual Evaluation Understudy)}.
BLEU measures the n-gram overlap between the generated text and the reference text, providing an estimate of lexical similarity and generation fluency. It is widely used for evaluating machine translation and other text generation tasks ~\cite{papineni2002bleu}.
% NK fix: missing full stop

\textbf{ROUGE.}
We report ROUGE-1 (also known as ROUGE-1-F1) and ROUGE-L (aka ROUGE-L-F1), which are commonly used for evaluating summarization and text generation systems. ROUGE-1 measures unigram overlap between the generated and reference texts, while ROUGE-L captures sequence similarity based on the longest common subsequence, reflecting structural alignment between sentences ~\cite{lin2004rouge}. For ROUGE, we report the F1 score, which balances precision and recall.

\textbf{BERTScore.}
To complement lexical overlap metrics, we also employ BERTScore, which evaluates semantic similarity using contextual embeddings derived from pretrained transformer models. This metric compares token embeddings between generated and reference texts to capture deeper semantic relationships. Following prior work, we report the F1 score as the main evaluation measure \cite{zhang2020bertscore}.

% NK change 1: pointer to the metrics of the RAG evaluation (reviewer comments 3, 6)
\textbf{RAG evaluation metrics.} The RAG-based KB is evaluated with a separate set of metrics: four metrics of the RAGAS framework~\cite{es2024ragas} (faithfulness, answer correctness, context precision and context recall), as well as the retrieval measures Hit@1 and Hit@20. Since these metrics assess answers to questions rather than generated abstracts, they are described together with the RAG evaluation in Subsection~\ref{sec:rag_eval}.

%\subsection{Experimental Setup}
\subsection{LM Studio Configurations}
%This subsection is dedicated to provide a list of configurations used for LM Studio.
To make repeating the experiments possible, all the LM Studio configurations were saved as "Presets". In LM Studio, a Preset refers to the way of bundling together a system prompt and other parameters into a single configuration.
However, the combinations of software and hardware  used to run our experiments are shown in  Appendix~\ref{app:sub_sw_hw}. % NK fixes (typos): "show" -> "shown"; "Apependix" -> "Appendix"
% NK change 2: the RAG evaluation used Ollama, not LM Studio (reviewer comments 3, 6)
The RAG evaluation does not use LM Studio. Rather, its models were run locally through Ollama, with the settings listed in Table~\ref{tab:rag_eval_setup}.

% \subsubsection{Software and Hardware} The necessary software for starting this experimental setup falls into three categories.

% \begin{itemize}
% \item	Software Packages: The LM Studio \textit{Version 0.3.31} or higher must be installed on hardware with sufficient computational power.

% \item	Local LLMs: We utilized three local LLMs that belong to different families.

% \item	Python Libraries: Creating a virtual environment using the Python built-in \textbf{venv} module is a prerequisite for this part. Being done with that, the needed Python libraries can be installed using the Python PIP, which is a package manager for installing the Python libraries. The following libraries were installed in the created virtual environment: 
% \end{itemize}

% \subsubsection{Hardware}
% %\paragraph{Hardware}
% The experiments were carried out on a Windows machine (MS Windows 11 installed) equipped with an NVIDIA RTX PRO 3000 Blackwell Generation laptop GPU.
\subsection{Model Evaluation and Preprocessing}
\label{subsec:eval}
%In the filed of NLP, generating abstracts is treated as a case of text summarization, and the later is treated as a case text generation.

%%To choose the right evaluation metrics, we need to determine the brand category of abstracts generation in the filed of NLP. In NLP, automatic abstract generation is typically modeled as a specialized text summarization task, and text summarization itself is widely categorized under natural language generation (text generation) within NLG/NLG-style sequence‑to‑sequence models \citep{becker2024textgenerationsystematicliterature}. 

Next, we will evaluate the different LLMs. We will report our findings as part of Section~\ref{sec:results}, using the aforementioned evaluation metrics. To do so, we iterate over the rows in the generated abstract based on \textbf{[\text{paper\_id}]}. The \textbf{[\text{gen\_abstract}]} column will be evaluated in comparison to their corresponding ground truth abstracts, i.e. the \textbf{[abstract]} column. Before stepping into the evaluation of the local LLMs using the aforementioned metrics, we noticed that a high-level preprocessing is necessary for the abstracts generated by the Meta-Llama model. For instance, the Llama-3.1-8B model was affected by the prompt. As a consequence, it provides long responses that start with a sentence like ``Here is a new formal academic abstract of exactly 120-130 words in one continuous paragraph:''. This noise was simply replaced with an empty string. % NK fixes (grammar): sentence fragment ". We" -> ", we"; ", As" -> ". As"; comma splice ", this" -> ". This"
% NK change 1: link this preprocessing to the RAG evaluation and quantify it from the indexed abstracts (reviewer comments 3, 5)
The same cleaned abstracts are indexed in the KB for the RAG evaluation (Subsection~\ref{sec:rag_eval}). Out of these 200 outputs, such an introductory sentence appeared 199 times; one output contained a refusal message instead of an abstract and was left unchanged.

% \begin{itemize}
% \item Bilingual Evaluation Understudy (BLEU):
% \item Recall-Oriented Understudy for Gisting Evaluation (ROUGE):
% \item BERTScore: This metric was calculated with \textit{rescale\_with\_baseline} set to False.
% \end{itemize}

%%%%%%%%%%%%%%%%%%%%%%%%%%%%%%%%%%%%%%%%%%
\section{Results}
\label{sec:results}

%\commentRed{Provide a concise and precise description of the experimental results, their interpretation as well as the experimental conclusions that can be drawn.} -- raw average

In this section, we present a quantitative evaluation of the local 
LLMs on the 200-article arXiv dataset and using the BLEU, ROUGE, and BERTScore metrics. % NK fixes: "evaluations" -> "evaluation" (grammar); dataset is not CS/ML only (reviewer comment 4)
 Our evaluation considers both long and short prompting strategies. It additionally shows the average raw and cleaned length for the ground truth and generated abstracts. % NK change 1: third part for the RAG evaluation results (reviewer comments 8, 14, 15)
We divided this into two dimensions: (i) a comparative study; and (ii) an ablation study. The results of the RAG-based KB evaluation are reported in Subsection~\ref{sec:rag_results}.

\subsection{Comparative Study}

This part was established using the default settings of LM Studio (see Figure~\ref{fig:lmstudio}) and divided into two steps. First, we evaluated a random subset of 50 articles in three iterations to validate the pipeline. Then, we reported the final results on the full 200-article corpus. The results obtained in the first step (i.e., using 50 articles) are shown in Table~\ref{tab:Iter-Results}, while the remaining results from step two (i.e., using 200 articles) are shown in Table~\ref{tab:allresults}. % NK fix (grammar): "Here in and after" -> "Hereafter"
Hereafter, we report the performance of the three local LLMs as calculated using the aforementioned evaluation metrics.  

% \begin{figure}[H]
% %\isPreprints{\centering}{} % Only used for preprints
% \includegraphics[width=6.0 cm]{img/lmstudio}
% \caption{The default settings for LM Studio.\label{fig:lmstudio}}
% \end{figure}

\begin{figure}[H]
\centering
\begin{minipage}{0.45\textwidth}
    \centering
    \includegraphics[width=\linewidth, height=5.5cm]{img/LM-StudioBasicSetting.png}
\end{minipage}
\hfill
\begin{minipage}{0.45\textwidth}
    \centering
    \includegraphics[width=\linewidth]{img/lmstudio} % second image
\end{minipage}
\caption{The default settings for LM Studio apply to all LLMs in the comparative study.}
\label{fig:lmstudio}
\end{figure}


% NK note (no change made): reviewer comment 9 asks for confidence intervals or significance tests for Tables 4-7 (e.g., BLEU differences between iterations).
\begin{table}[H]
%\small
\footnotesize
\caption{Performance comparison of article summarization using small LLMs across prompts and iterations,  reported in Mean ± SD format. Each model was evaluated in three iterations using the same subset of 50 articles.\label{tab:Iter-Results}}
%\isPreprints{\centering}{% This command is only used for ``preprints''.
	\begin{adjustwidth}{-\extralength}{0cm}
%} % If the paper is ``preprints'', please uncomment this parenthesis.
%\isPreprints{\begin{tabularx}{\textwidth}{CCCC}}{% This command is only used for ``preprints''.
		\begin{tabularx}{\fulllength}{lCCrcccc}
%} % If the paper is ``preprints'', please uncomment this parenthesis.
			\toprule
			\textbf{Model}	& \textbf{Prompt}	& \textbf{Iter.\textsuperscript{*}}	& \textbf{Avg. W. Len}     & \textbf{BLEU} & \textbf{ROUGE-1} & \textbf{ROUGE-L} & \textbf{BERTScore}\\

\midrule
\multirow{3}{*}{Mixtral-8$\times$7B} &
\multirow{3}{*}{Long} &
    1 & 230 & 5.18 ± 9.15   & 30.50 ± 12.84   & 17.89 ± 9.10   & 82.98 ± 3.24   \\
& & 2 & 226  & 3.87 ± 4.80   & 30.35 ± 10.32   & 17.84 ± 7.03   & 82.98 ± 2.79   \\
& & 3 & 228  & 6.99 ± 10.08   & 35.91 ± 12.97   & 21.41 ± 11.19   & 84.57 ± 2.81   \\

\midrule
\multirow{3}{*}{Qwen3-VL-8B} &
\multirow{3}{*}{Short} &
    1 & 159    & 3.51 ± 2.46   & 34.65 ± 7.08   & 17.08 ± 4.10   & 84.80 ± 1.32 \\
& & 2 & 171     & 3.42 ± 2.81   & 34.26 ± 7.76   & 17.01 ± 4.79   & 84.72 ± 1.39   \\
& & 3 & 162     & 3.53 ± 2.83   & 33.17 ± 7.87   & 16.49 ± 4.51   & 84.54 ± 1.67   \\

\midrule
\multirow{3}{*}{Llama-3.1-8B} &
\multirow{3}{*}{Short} &
%Length: 160, 149, 147
    1 & 163 & 9.07 ± 5.62   & 43.16 ± 9.45   & 23.83 ± 7.13   & 86.68 ± 1.73   \\
& & 2 & 138  & 10.69 ± 8.70   & 44.44 ± 10.92   & 26.40 ± 10.33   & 87.07 ± 2.32   \\
& & 3 & 137  & \textbf{12.48 ± 8.59}   & \textbf{45.91 ± 10.31}   & \textbf{27.34 ± 9.98}   & \textbf{87.26 ± 2.27}   \\

                               
			\bottomrule
		\end{tabularx}
%		\isPreprints{}{% This command is only used for ``preprints''.
	\end{adjustwidth}
%} % If the paper is ``preprints'', please uncomment this parenthesis.
	\noindent{\footnotesize{* \textbf{Iter.} is an abbreviation for Iteration.}}
\end{table}


% NK note (no change made): the numbers in this paragraph do not match Table Iter-Results. Text: BLEU 6.74 -> 4.94 -> 9.21; table: 5.18 -> 3.87 -> 6.99. Text: Llama ROUGE-L 45.07 and BERTScore 86.70; table: ROUGE-L 23.83-27.34 (ROUGE-1 is 43.16-45.91) and BERTScore 86.68-87.26. Authors of this part: please check which values are correct.

Table~\ref{tab:Iter-Results} compares the performance of the three local LLMs (Mixtral-8$\times$7B, Qwen3-VL-8B, and Llama-3.1-8B) in article summarization across different prompts and iterations.
%%Given that Mixtral-8$\times$7B model cannot be instructed using a "Short" prompt. 
The reported results show clear differences in summarization quality across the three evaluated models.
Among the three models, Llama-3.1-8B achieves the best summarization performance, with consistent improvements in BLEU, ROUGE-1, ROUGE-L, and BERTScore across all iterations.
Although Qwen3-VL-8B scores decrease slightly (except BLEU score) in the third iteration, it shows limited variation across iterations.
Mixtral-8$\times$7B had modest improvement overall; it mainly improves slightly from iteration 2 to iteration 3.
Additionally, we used the paired t-test to calculate the statistical significance of the differences between the three models in each iteration.
In iterations 1 and 3, Llama-3.1-8B outperformed both Qwen3-VL-8B and Mixtral-8$\times$7B across all metrics. All the results are statistically significant. However, iteration 2 showed no significant results. %%Within the model itself, 

Furthermore, the average length of generated abstracts also varied noticeably between the three models, with Mixtral-8$\times$7B generating substantially longer abstracts (approximately 230 words) than the other two models (approximately 145 - 165 words). Overall, the shorter abstracts \textbf{(137 words)} generated by Llama-3.1-8B were associated with higher automatic evaluation scores, closer to the  average length of gold-standard abstracts (125 words). However, the longer outputs (230 words) from Mixtral-8$\times$7B did not consistently lead to better performance.

%DEL:% The Mixtral-8$\times$7B model shows a notable increase in BLEU score from 6.74 to 9.21 over three iterations. 
%DEL:% Initially after the first iteration, the Bleu score drops from 6.74 to 4.94 and then again after the third iteration it goes up to 9.21. Llama-3.1-8B achieves the highest ROUGE-L scores of 45.07 and BERTScore of 86.70 across all prompts and iterations. These results suggest its superior summarization capabilities compared to Qwen3-VL-8B. """Furthermore, the average length of generated abstracts also varied noticeably between the three models, with Mixtral-8$\times$7B generating substantially longer abstracts (approximately 230 words) than the other two models (approximately 150 - 160 words). Overall, the shorter abstracts \textbf{(150 words)} generated by Llama-3.1-8B were associated with higher automatic evaluation scores, closer to the  average length of gold-standard abstracts (125 words), while longer outputs (230 words) from Mixtral-8$\times$7B did not consistently translate into better performance""". 


%DEL:%\commentVP{We performed paired t-tests to check whether the results obtained were statistically significant or not. For comparison within the iterations, given that the Mixtral-8$\times$7B model cannot be instructed using a "Short" prompt, Mixtral-8$\times$7B had modest improvement overall; it mainly improves from iteration 2 to iteration 3 (BLEU p = 0.014). Qwen3-VL-8B had no meaningful improvement, whereas Llama-3.1-8B improves significantly from iteration 1 to iteration 3 (BLEU p = 0.0033, ROUGE-1 p = 0.0109). This improvement occurs while the average generated length decreases from 163 to 137 words. Iteration 2 showed no significant gain.}

%DEL:%\commentVP{When comparing across models for iteration 1, Llama-3.1-8B outperforms both Mixtral-8$\times$7B (all metrics p < 0.02) and Qwen3-VL-8B (all metrics p < 0.0001). Mixed results were observed when we compared Mixtral-8$\times$7B vs Qwen3-VL-8B. Qwen3-VL-8B was better on ROUGE-1(p=0.012) and BERTScore (p < 0.0001), but BLEU (p = 0.20) and ROUGE-L (p = 0.57) were not statistically significant. Finally, for iteration 3, Llama-3.1-8B still outperforms and is statistically significant as well, since Mixtral-8$\times$7B (all metrics p < 0.01) and Qwen3-VL-8B (all metrics p < 0.0001). Overall, Llama-3.1-8B with the Short prompt provides the strongest results among the evaluated configurations.}


Similar to Table~\ref{tab:Iter-Results}, Table~\ref{tab:allresults} compares the performance of the three local LLMs but it uses the entire sub-dataset of 200 articles.
Table~\ref{tab:allresults} demonstrates clear performance differences among the three evaluated LLMs for abstract generation. 
Llama-3.1-8B achieves the best summarization performance among the three models under the tested configurations.
It also produces the shortest generated abstracts on average, with an average length of 138 words.
Qwen3-VL-8B provides competitive semantic similarity scores, but remains behind Llama-3.1-8B model in overall summarization quality. The Mixtral-8$\times$7B model generates substantially longer abstracts, and does not adhere to the given prompt. 
Consequently, we consider it outside the competition. Overall, the results suggest that higher abstract length alone does not necessarily lead to improved quality.
We calculated the $p$-value ($p$) to check the superiority of Llama-3.1-8B model compared to both models in Table~\ref{tab:allresults}. 
We found that its reported results are statistically significant in both lexical overlap/coverage and semantic similarity.
%DEL%We found that in both lexical coverage (lexical overlap) and semantic similarity. 
Specifically, the $p$-values for BLEU, ROUGE-1, ROUGE-L, and BERTScore are all less than 0.0001.
Table~\ref{tab:p_value_comparisons} shows all the $p$-values that relate to model comparison.

%DEL:%\commentVP{For Table 6, Llama-3.1-8B is highly statistically significant across all the metrics compared to both models. For the full dataset, compared to Mixtral-8$\times$7B, it is significantly superior. Here are the p-values for its BLEU (p < 0.0001), ROUGE-1 (p < 0.0001), ROUGE-L (p < 0.0001), and BERTScore (p < 0.0001), and the same p-values hold for Qwen3-VL-8B as well. For lexical overlap and semantic similarity, results show an advantage for Llama-3.1-8B.  Overall, For Mixtral-8$\times$7B vs Qwen3-VL-8B, Mixtral-8$\times$7B is slightly better overall. Significantly in BLEU (p = 0.0428) and BERTScore (p < 0.0001), but ROUGE-L is not statistically significant since its p-value is higher than the standard threshold.}

%DEL:Old Text%Llama-3.1-8B achieves the highest overall results, scoring highest on all metrics and producing the shortest abstracts. Qwen3-VL-8B shows moderate performance for all metrics except for the BLEU and slightly in ROUGE-L, where it is the lowest. "Qwen3-VL-8B provides competitive semantic similarity scores but remains behind Llama-3.1-8B model in overall summarization quality. The Mixtral-8$\times$7B model generates substantially longer abstracts, and does not adhere to the given prompt. Consequently, we consider it out of the competition." Overall, the results suggest that higher abstract length alone does not necessarily lead to improved quality. Llama-3.1-8B offers the best balance between concise abstracts generation and strong evaluation performance.  

%\begin{listing}[H]
%\caption{Title of the listing}
%\rule{\columnwidth}{1pt}
%\raggedright Text of the listing. In font size footnotesize, small, or normalsize. Preferred format: left aligned and single spaced. Preferred border format: top border line and bottom border line.
%\rule{\columnwidth}{1pt}
%\end{listing}

% \begin{table}[H] 
% \caption{The overall performance of article summarization using small LLMs, evaluated on the full dataset of 200 articles and reported in Mean ± SD format.\label{tab:allresults}}
% \begin{tabularx}{\textwidth}{lrrrrr}
% \toprule
% \textbf{Model}	& \textbf{Avg. W. Len}     & \textbf{BLEU} & \textbf{ROUGE-1} & \textbf{ROUGE-L} & \textbf{BERTScore}\\
% \midrule
% Mixtral-8$\times$7B		& 230		& 5.12  & 30.14 & 17.62 & 82.85 \\
% Qwen3-VL-8B		& 169			& 3.66  & 34.51 & 17.10 & 84.77 \\
% Llama-3.1-8B		& 138		& 12.96 & 45.14 & 26.17 & 86.82 \\%Length:147
% \bottomrule
% \end{tabularx}
% \end{table}

%%A small change in the reported BLEU was noticed.
\begin{table}[H] 
\small
\caption{The overall performance of article summarization using small LLMs, evaluated on the full dataset of 200 articles and reported in Mean ± SD format.\label{tab:allresults}}
%\isPreprints{\centering}{} % Only used for preprints
%\begin{tabularx}{\textwidth}{CCCccc}
\begin{tabularx}{\textwidth}{lrcccc}
\toprule
\textbf{Model}	& \textbf{Avg. W. Len}     & \textbf{BLEU} & \textbf{ROUGE-1} & \textbf{ROUGE-L} & \textbf{BERTScore}\\
\midrule
Mixtral-8$\times$7B		& 230		& 4.13 ± 6.75  & 30.14 ± 11.33 & 17.62 ± 8.10 & 82.85 ± 2.73 \\
Qwen3-VL-8B		    & 169	    & 3.20 ± 2.49  & 34.51 ± 7.23 & 17.10 ± 4.12 & 84.77 ± 1.32 \\
Llama-3.1-8B		& 138		& \textbf{11.50 ± 8.76} & \textbf{45.94 ± 10.83} & \textbf{27.16 ± 10.20} & \textbf{87.38 ± 2.36} \\%Length:147
\bottomrule
\end{tabularx}
%\noindent{\footnotesize{\textsuperscript{1} Tables may have a footer.}}
\end{table}


\begin{table}[htbp]
\small
\centering
\caption{Statistical significance ($p$-values) of performance differences between models across the full dataset.}
\label{tab:p_value_comparisons}
\resizebox{\textwidth}{!}{%
\begin{tabular}{lcccc}
\toprule
\textbf{Model Comparison} & \textbf{BLEU} & \textbf{ROUGE-1} & \textbf{ROUGE-L} & \textbf{BERTScore} \\ 
\midrule
Meta-Llama-3.1-8B vs. Mixtral-8x7B & $<0.001^*$ & $<0.001^*$ & $<0.001^*$ & $<0.001^*$ \\ 
Meta-Llama-3.1-8B vs. Qwen3-VL-8B  & $<0.001^*$ & $<0.001^*$ & $<0.001^*$ & $<0.001^*$ \\ 
Mixtral-8x7B vs. Qwen3-VL-8B       & $0.0428^*$ & $<0.001^*$ & $0.3404$   & $<0.001^*$ \\ 
\bottomrule
\multicolumn{5}{l}{\footnotesize  $^*$Statistically significant ($p < 0.05$).} \\
\end{tabular}%
}
\end{table}

Furthermore, our comparison went beyond comparing the reported metrics to include: (i) word-cloud visualization; and (ii) BERTopic visualization. % NK fix (typo): "world-cloud" -> "word-cloud" (here and in the caption below)
Figure~\ref{fig:WordCloud} shows the word-cloud representation of the ground truth (gold standard) abstracts and the generated abstracts using Llama-3.1-8B model. Some domain-specific technical terms can be observed  in the ground truth, such as: "network”, "problem", "time", "field", "equation", "algorithm" ... etc. 
In the AI-generated, some terms overlap such as "network", "time", "field", "system". 
However, it shifts towards more generic usage of terms rather than technical ones, such as "implications", "understanding", "investigate" ... etc. 
Overall, the AI-generated content is good for academic purposes, but lacks a bit of information density for that specific domain.

% Example of a figure that spans the whole page width and with subfigures. The same concept works for tables, too.
\begin{figure}[H]
%\isPreprints{}{% This command is only used for ``preprints''.
\begin{adjustwidth}{-\extralength}{0cm}
\centering
%} % If the paper is ``preprints'', please uncomment this parenthesis.
\subfloat[\centering]{\includegraphics[width=8.5cm]{img/arxiv_abstracts_wc.png}}
\hfill
\subfloat[\centering]{\includegraphics[width=8.5cm]{img/llama3_abstracts_wc.png}}\\

%%22-17-2026 OH:
% \subfloat[\centering]{\includegraphics[width=8.70cm]{img/nltk_arxiv_word_freq}}
% \hfill
%%22-12-2025 OH:
% \subfloat[\centering]{\includegraphics[width=8.70cm]{img/nltk_llama_word_freq}}
%\isPreprints{}{% This command is only used for ``preprints''.
\end{adjustwidth}
%} % If the paper is ``preprints'', please uncomment this parenthesis.
\caption{(\textbf{a}) The word-cloud of the 200 gold standard abstracts. (\textbf{b}) The word-cloud of the 200 generated abstracts using Llama-3.1-8B model. \label{fig:WordCloud}}
\end{figure} 


% NK note (no change made): this paragraph says the clusters were identified by the NLTK library (and the file is nltk_word_frequency_vertical.pdf), but the caption and the list above say BERTopic. Authors of this part: please make them consistent.
Figure~\ref{fig:BERTopic} shows bar charts that compare the top 20 word clusters identified by the NLTK library in the gold standard dataset with the Llama-3.1-8b generated abstracts. 
%The horizontal bar chart demonstrates the implementation of NLTK library for comparison of the two datasets. 
The data has been lemmatized, and standard English stop-words have been removed. If observed, we see that gold standard dataset uses more of the nouns such as "network”, "model", "time", "algorithm" ... etc. On the other hand, in AI-generated abstracts, there are common words to the gold standard dataset, but the use of verbs is more dominant such as "approach", "using", "study", etc. Hence, the AI-generated words are more generic as compared to gold standard dataset.

\begin{figure*}[ht]
\centering
\includegraphics[width=0.99\textwidth, height=0.52\textheight]{img/nltk_word_frequency_vertical.pdf}
\caption{The top word embedded clusters found by BERTopic.}
\label{fig:BERTopic}
\end{figure*}

%TODO, 2026-07-31 OH:
%% If the documentclass option "submit" is chosen, please insert a blank line before and after any math environment (equation and eqnarray environments). This ensures correct linenumbering. The blank line should be removed when the documentclass option is changed to "accept" because the text following an equation should not be a new paragraph.

%% Example of a page in landscape format (with table and table footnote).
%\startlandscape
%\begin{table}[H] %% Table in wide page
%%\isPreprints{\centering}{} % This command is only used for ``preprints''.
%\caption{This is a very wide table.\label{tab3}}
%	\begin{tabularx}{\textwidth}{CCCC}
%		\toprule
%		\textbf{Title 1}	& \textbf{Title 2}	& \textbf{Title 3}	& \textbf{Title 4}\\
%		\midrule
%		Entry 1		& Data			& Data			& This cell has some longer content that runs over two lines.\\
%		Entry 2		& Data			& Data			& Data\textsuperscript{1}\\
%		\bottomrule
%	\end{tabularx}
%%\isPreprints{}{% This command is only used for ``preprints''.
%	\begin{adjustwidth}{+\extralength}{0cm}
%%} % If the paper is ``preprints'', please uncomment this parenthesis.
%		\noindent\footnotesize{\textsuperscript{1} This is a table footnote.}
%%\isPreprints{}{% This command is only used for ``preprints''.
%	\end{adjustwidth}
%%} % If the paper is ``preprints'', please uncomment this parenthesis.
%\end{table}
%\finishlandscape

%% 22-12-2025 OH:
% Please punctuate equations as regular text. Theorem-type environments (including propositions, lemmas, corollaries etc.) can be formatted as follows:
% %% Example of a theorem:
% \begin{Theorem}
% Example text of a theorem.
% \end{Theorem}

%% 22-12-2025 OH:
%% Example of a proof:
% \begin{proof}[Proof of Theorem 1]
% Text of the proof. Note that the phrase ``of Theorem 1'' is optional if it is clear which theorem is being referred to.
% \end{proof}
\subsection{Ablation Study}

We have carried out several standards evaluation experiments so far. Based on that, it can be concluded that \textit{Llama-3.1-8B} outperformed the other two LLMs. Now is the time to conduct an ablation study. Such a study involves systematically removing or altering individual parameters of the LLM. For the scope of this paper, we change the prompt design, and tune the decoding or sampling parameters of the \textit{Llama-3.1-8B} model in two separate experiments. In our settings, we start by studying the effect of prompt design, and then focus on decoding parameters.

\subsubsection{Prompt Design}
\label{sub:prompt_design}

As inspired by \citet{zhang-etal-2024-benchmarking}, our experimental settings are based on the zero-shot prompting. However, we changed the prompt design in terms of prompt structure and constraints. Table~\ref{tab:prompts-design}, in Appendix~\ref{app:sub_prompt_des}, summarizes the prompts  used for abstract generation in Llama-3.1-8B LLM.  As explained below, each prompt configuration incrementally adds specific functional constraints.
%TODO: OH% Table~\ref{tab:prompt_ablation_matrix} lists the incremental prompt ablation component matrix that illustrates the exact parameter changes between P1 and P5.

\begin{itemize}
    \item P1 (Plain): An unconstrained prompt that requests a formal abstract without length, format, or content rules.
    \item P2 (Semi-structured with paragraph formatting): This prompt adds a structural constraint that enforces output generation as one continuous paragraph.
    \item P3 (Semi-structured with length constraint): This prompt adds a quantitative length constraint of 120–130 words.
    \item P4 (Structured with length and content constraints): This prompt adds content coverage directives, which mandate the explicit inclusion of the four standard components of an academic abstract : Background, Methods, Results, and Conclusions.
    \item P5 (Fully structured with LSCSR constraints): This prompt introduces the LSCSR (Length + Style + Content + Strict Rules) constraints. It also adds negative constraints ("Do NOT include feedback, suggestions, critiques, or numbered lists") and strict single-output directives ("Write ONLY a new formal academic abstract...") to eliminate LLM meta-chatter and preamble noise.  
\end{itemize}

Table~\ref{tab:prompt_ablation} reports the results of evaluating the different prompt designs for abstract generation using Llama-3.1-8B model with 50 and 100 articles respectively. Prompt design has a clear impact on abstract generation quality across both evaluation settings. 
The prompt design P1 was set as the baseline and compared with the prompt designs P2 - P5.
We observed that introducing only a semi-structured paragraph formatting (P2) gives a highly significant improvement for both 50 articles and 100 articles.
To test whether the reported results are statistically significant for P1 vs. P2, we calculated the $p$-values for both cases: 50 and 100 articles.
We found that the reported results are statistically significant ($p < 0.05$) for all metrics.

%DRL:%\commentVP{The $p$-values for the metrics were reported as BLEU (p = 0.043), ROUGE-1 (p = 0.001), ROUGE-L (p = 0.025), and BERTScore (p = 0.0031) for 50 artica les and BLEU (p = 0.0048), ROUGE-1 (p = 0.0012), ROUGE-L (p = 0.0002), and BERTScore (p = 0.00018) for 100 articles.} 
Moving from a plain prompt to a semi-structured prompt consistently improves the overall performance. This means that the model benefits from providing explicit guidance. Similar performance measures were observed for both 50 and 100 articles. Overall, the main difference is in the Avg. W. Len observed. For a prompt with plain structure, it generates about 200 words as compared to the fully structured prompt words which generates nearly 150.
Using a fully structured prompt (P5) reduces the Avg. W. Len, while maintaining competitive ROUGE and BERTScore results. 
Other than that, all configurations (prompt structures) produce comparable results across the evaluation metrics. 

%DEL%\commentVP{Additionally, the tests also confirm that length, content, and style constraints can hinder the LLM. Prompts P3, P4, and P5 failed to yield statistically significant results since the p-values are greater than the standard threshold.} 
This means that no significant improvements were observed in any of the prompt structures. Thus, the prompt design has only a limited effect on the overall summarization performance when using the Llama-3.1-8B model.

P2 scored highest on all four metrics for both sample sizes. P5 scored slightly lower, for example BLEU 11.42 vs.\ 13.91 for 50 articles and 10.53 vs.\ 13.37 for 100, while BERTScore differed by less than 0.6 points. These differences are small compared with the standard deviations, which are large for every prompt. P5, however, produced the shortest abstracts (154 and 150 words), the closest to the 120--130-word target. Since P5 gave more control over length, style and content, and followed a strict set of rules, we decided to proceed with the fully structured prompt for the rest of our cases, accepting a small loss in lexical overlap.

%%\begin{adjustwidth}{-\extralength}{0cm}
%Overfull \hbox (15.63217pt too wide) in paragraph at lines 838--850


\begin{table}[H]
\centering
\caption{Ablation Study on Prompt Design: Evaluation summary for the abstract generation using Llama-3.1-8B, tested using 50 and 100 articles, and reported in Mean ± SD format.}
\label{tab:prompt_ablation}

\resizebox{\textwidth}{!}{

\begin{tabular}{lllrrcccccccc}
\toprule
\textbf{ID} & \textbf{Structure} & \textbf{Constraints}
& \multicolumn{2}{c}{\textbf{Avg. W. Len}}
& \multicolumn{2}{c}{\textbf{BLEU}}
& \multicolumn{2}{c}{\textbf{ROUGE-1}}
& \multicolumn{2}{c}{\textbf{ROUGE-L}}
& \multicolumn{2}{c}{\textbf{BERTScore}} \\
\cmidrule(lr){4-5}
\cmidrule(lr){6-7}
\cmidrule(lr){8-9}
\cmidrule(lr){10-11}
\cmidrule(lr){12-13}
 &  & 
& \textbf{50} & \textbf{100}
& \textbf{50} & \textbf{100}
& \textbf{50} & \textbf{100}
& \textbf{50} & \textbf{100}
& \textbf{50} & \textbf{100} \\
\midrule
P1 & Plain         & None                  & 204 & 199 & 11.37 ± 5.85 & 10.77 ± 6.99 & 45.03 ± 9.91  & 44.10 ± 11.16 & 26.43 ± 6.75  & 25.64 ± 7.61  & 87.07 ± 1.62 & 86.76 ± 1.93 \\
P2 & Semi-struct.  & Paragraph format      & 161 & 171 & \textbf{13.91 ± 8.52 }& \textbf{13.37 ± 9.34} & \textbf{48.95 ± 10.39} & \textbf{46.92 ± 11.92} &  \textbf{29.31 ± 10.43} & \textbf{28.76 ± 10.33}  & \textbf{87.76 ± 2.17} & \textbf{87.49 ± 2.35 }\\
P3 & Semi-struct.  & Length (120--130 W.)  & 156 & 154 & 11.68 ± 8.10 & 11.76 ± 8.12 & 45.65 ± 10.00 & 45.63 ± 11.51 & 26.90 ± 10.01 & 26.97 ± 10.05 & 87.18 ± 2.10 & 87.14 ± 2.31 \\
P4 & Structured    & Length + Content      & 161 & 155 & 11.50 ± 8.17 & 11.56 ± 7.61 & 46.18 ± 10.16 & 45.21 ± 11.44 & 26.80 ± 9.40  & 26.05 ± 9.08  & 87.17 ± 2.31 & 87.14 ± 2.38 \\
P5 & Fully struct. & LSCSR*                & 154 & 150 & 11.42 ± 7.49 & 10.53 ± 8.54 & 45.53 ± 10.63 & 44.90 ± 11.28 & 27.15 ± 9.25  & 26.44 ± 9.50  & 87.23 ± 2.07 & 87.04 ± 2.39 \\
\bottomrule
\end{tabular}

}
\noindent{\footnotesize{\textbf{*LSCSR}: Length + Style + Content + Strict Rules.}}
\end{table}


\subsubsection{Decoding Parameters}
\label{subs:dec_para}

Following the work by \citet{Hamed2026IC}, we consider three decoding parameters, particularly temperature (denoted as temp), \text{top-p} (aka nucleus sampling), and \text{top-k} sampling. 
The following text describes these parameters in LLMs context based on the work provided by ~\citet{nguyen2025turningheatminpsampling,chang2023kldivergenceguidedtemperaturesampling}.

\begin{itemize}
    \item Temperature or temp: This setting controls the level of randomness with which a model selects words during text generation. In our case, it refers to the generation of abstract text. Thus, it determines the degree of freedom, which the model has, when selecting the next word in a sequence of words that represent the generated abstract text.
    \item Top-p or top\_p (aka nucleus sampling): This setting determines how many possible words the model considers before making a selection. Thus, it filters the potential word pool based on a cumulative probability threshold.
    \item Top-k or top\_k sampling: This setting represents another sampling parameter that used to control word selection. While top-p selects a dynamic number of words based on a cumulative probability threshold, the top-k limits the model to a fixed number of the most probable subsequent words.
\end{itemize}

As per the "Prompt Design" experiments, we noticed that the number of random articles had a minimal effect on the reported metrics for 50 and 100 articles. 
Consequently, we carried out the "Decoding Parameters" experiments on the same 50 articles as in Subsection~\ref{sub:prompt_design},  the prompt design with the ID\#~P5, as well as the LM Studio configurations (aka Presets).
These configurations are much stronger because they isolates a single decoding parameter in each experiment. Table~\ref{tab:decodingresults} summarizes the results obtained from this joint setting.

\begin{table}[H] 
%\small % Change table font size
\caption{The effect of decoding parameters on article summarization performance using Llama-3.1-8B, assessed using the same 50 random articles and reported in Mean ± SD format.\label{tab:decodingresults}}
\resizebox{\textwidth}{!}{
%\begin{tabularx}{\textwidth}{lrrrrr}
\begin{tabular}{llrcccc}
\toprule
\textbf{ID} & \textbf{Configuration}	& \textbf{Avg. W. Len}     & \textbf{BLEU} & \textbf{ROUGE-1} & \textbf{ROUGE-L} & \textbf{BERTScore}\\
\midrule
C1 & Temp=0.30, \text{top\_p}=0.85, \text{top\_k}=30                & 147 & 10.13 ± 8.55 & 44.79 ± 11.11 & 26.25 ± 10.25 & 87.03 ± 2.34 \\
C2 & Temp=0.65, \text{top\_p}=0.80, \text{top\_k}=30	            & 141 & 10.85 ± 9.93 & 45.60 ± 11.25 & 26.53 ± 11.22 & 87.25 ± 2.48 \\
C3 & Temp=0.65, \text{top\_p}=0.85, \text{top\_k}=30	(reference) & 151 & 10.94 ± 8.20 & 45.25 ± 9.96 & \textbf{27.15 ± 9.88} & 87.30 ± 2.20 \\%\textsuperscript{1}
C4 & Temp=0.65, \text{top\_p}=0.90, \text{top\_k}=30		        & 155 & 10.79 ± 8.13 & 45.28 ± 10.07 & 26.80 ± 10.07 & 87.11 ± 2.28 \\
C5 & Temp=1.00, \text{top\_p}=0.85, \text{top\_k}=30		        & 152 & \textbf{10.97 ± 8.12} & \textbf{46.57 ± 10.78} & 26.51 ± 9.29 & \textbf{87.31 ± 2.29} \\
\bottomrule
\end{tabular}
}
%\noindent{\footnotesize{\textsuperscript{1} Tables may have a footer.}}
\end{table}

As per the results in Table \ref{tab:decodingresults}, the third configuration (C3) represents the baseline.
All configurations produce comparable results across the various evaluation metrics. 
C3 was then compared with the four other configurations: C1, C2, C4, and C5. 
All the resulting $p$-values were greater than the standard threshold value. 
This indicates that none of the results from decoding parameters adjustments led to statistically significant  performance. Thus, the decoding parameters have only a limited effect on the overall summarization performance using the Llama-3.1-8B model. The reference configuration (Temp=0.65, top\_p=0.85, top\_k=30) provides the most balanced performance. It achieves the highest ROUGE-L score while maintaining competitive BLEU, ROUGE-1, and BERTScore values. 
Increasing the temperature from Temp $=$ 0.65 to Temp $=$ 1.00, slightly decreases the ROUGE-L. 
However, it improves the BLEU and ROUGE-1 scores; the BERTscore almost remains stable. 
The average word length (Avg. W. Len) remains relatively stable across different configurations. This indicates that the decoding parameters have minimal impact on output length. In conclusion, these results suggest that the reference configuration offers the best balance between lexical overlap and semantic similarity for this task. 
A key finding  from this statistical test is that Llama-3.1-8B’s summarization capabilities remain robust across a wide range of hyperparameters for this task.

%%%%%%%%%%%%%%%%%%%%%%%%%%%%%%%%%%%%%%%%%%
% \section{Discussion}
% \label{sec:discussion}

% Authors should discuss the results and how they can be interpreted from the perspective of previous studies and of the working hypotheses. The findings and their implications should be discussed in the broadest context possible. Future research directions may also be highlighted.

%%%%%%%%%%%%%%%%%%%%%%%%%%%%%%%%%%%%%%%%%%
\section{RAG-based Knowledge Base System}
\label{sec:KB}

As a proof of concept, we developed a RAG-based KB system (RAG-KBS), in which we index the documents using only their abstracts. This  is a very common approach \cite{aytar2024retrieval} to making our search system (i.e., the KB) faster and more accurate. % NK fix (grammar): "faster and accurate" -> "faster and more accurate"

\subsection{System Description}
Developing this RAG-KBS involves two main steps: (i) a RAG pipeline, which involves indexing and knowledge preparation (offline); and (ii) retrieval and generation (online), when end users search for the relevant papers in the VDB from the web interface.

%\subsection{RAG Pipeline} \subsection{Retrieval and Generation}
We utilized the LangChain \cite{chase2022langchain} Python library to implement a typical RAG pipeline comprising  the following sub-steps. Figure~\ref{fig:RAG} depicts these sub-steps.

\begin{itemize}
    \item Load Documents and Text Preprocessing: First, the system loads the external knowledge source, which in our case is a CSV file containing the abstracts generated by the Llama-3.1-8B model. 
    %\commentGN{@NK (Add the process)}: \commentVP{We load the file using \commentRed{Langchain's} CSVLoader keeping abstract as the main column. Additionally, we preprocess the text to \commentRed{remove the latex format backslashes.}}
    Before indexing, the introductory sentences that the model added to its output (e.g., ``Here is a combined abstract\ldots'') were removed (Subsection~\ref{subsec:eval}).

    \item Chunk the Documents: 
    The large documents (i.e. abstracts) are typically split into smaller sections called the "chunks". 
    Usually, this sub-step employs overlapping chunks because some important information may exist near chunk boundaries. 
    %%\ifshowchanges{\color{red}\textbf{[NK: deleted from here]}} \commentVP{We use \commentRed{recursive character text splitter} with a chunk size of \commentRed{800 (123 - 160 words)} to capture atleast 2 to 4 complete statements and overlap step of \commentRed{150} characters for keeping the consistency in the context across the characters for relation.} {\color{red}\textbf{[NK: to here]}}\fi 
    % NK change 1: the evaluated contexts are whole abstracts. In the RAGAS input, all 1,600 retrieved contexts (80 questions x 20) are complete abstracts, and 1,555 of them are longer than 800 characters, so no 800/150 splitting was applied in the evaluated configuration. VP: if the web system uses a splitter, we can mention it as a deployment option, but not for the evaluation.
    In our implementation, however, the abstracts are short (106--233 words, mean 137.6), so each abstract is stored as a single chunk: no splitting and no overlap are applied.

    \item Embed the Chunks: After splitting the abstracts into smaller chunks, the system must convert each chunk into a format that the LLM can understand mathematically. 
    Therefore, we embed those using the \textbf{nomic-embed-text} embedding model.
    % NK note (no change made): the embedding model of the KB is not in the evaluation files or code I have, so I cannot confirm bge-small-en-v1.5. VP: please confirm. (The RAGAS evaluation itself used nomic-embed-text, Table rag_eval_setup.)
    
    \item Store in a Vector Database and Indexing: As defined in \cite{brousseau2025llms}, vector databases are a special type of database designed specifically for storing vectors and the metadata that describes them. This makes them highly effective for storing embeddings. Thus, they are unlike a relational database management system (RDBMS) that operates search using keywords matching. An LLM needs the VDBs to employ semantic search (search by meaning). Furthermore, "they make RAG so much easier to implement" \cite{brousseau2025llms}.  That's why we need to use a VDB with the RAG-KBS. \citet{brousseau2025llms} described several VDBs, including Pinecone, Milvus, Redis and Elastic, FAISS (Facebook AI Similarity Search), USearch, and Elasticsearch. For this work, however, we stored and indexed it locally using ChromaDB because it operates locally with a lightweight setup and fully embedding  without any network call or cloud dependency. % NK note (no change made): ChromaDB cannot be checked from the evaluation files. VP: please confirm.
    
    \item Embed the User Query: This should be done using the same embedding model in the previous sub-step (i.e., embed the chunks). Thus, when an end user asks a question or searches for a document, the system must also convert the query into a vector embedding with the same formatting.
    \item Retrieve Relevant Chunks: ChromaDB stored the vector representation of every chunk, which in our case is one abstract per chunk. The user's query was converted into a vector by the same embedding model, and a dense similarity search retrieved the 20 most similar abstracts from the VDB, which were passed, together with their paper identifiers, to the LLM.
    %\commentGN{THIS STATEMENT FROM HERE WILL BE DELETED!!} 
    %%\ifshowchanges{\color{red}\textbf{[NK: deleted from here]}} \commentRed{and with weighted hybrid approach of sparse BM25 keyword search + dense vector search (sematics) for qwen2.5:7b due to computational restrictions}. {\color{red}\textbf{[NK: to here]}}\fi
  
    % NK fixes: the query is embedded by the embedding model, not the LLM (see the previous item); "retrieve" -> "retrieved". NK change 2: the evaluation used one configuration, so "for qwen2.5:14b" (Qwen2.5-14B is the RAGAS judge, not part of retrieval) and the hybrid sentence (as GN asked) are deleted; every question has 20 retrieved contexts in the RAGAS input, so "top 3" -> "top 20". NK note for VP: please confirm that the reranker was used in the evaluated run.
    
    \item Prompt Augmentation: This step is typically based on prompt injection. The retrieved context will be injected into the prompt template, and a new prompt will be augmented and fed to the LLM.
    
    \item LLM Response: Finally, the LLM generates a grounded and accurate response, which usually includes a list of documents sorted by their relevance. % NK note (no change made): the evaluation shows that answers are not always accurate (mean answer correctness 0.689 on valid questions, Table rag_scores); "grounded and accurate" could be softened to "grounded".
\end{itemize}

\begin{figure*}[ht]
\centering
\includegraphics[width=0.99\textwidth]{img/RAG_Updated.png}
\caption{The sub-steps of the RAG pipeline.}
\label{fig:RAG}
\end{figure*}


\subsection{Question Generation}
Before diving into system evaluation using RAGAS, we need to generate the questions and answers (Q\&A) dataset. We randomly selected 80 abstracts out of 200. For each abstract, Meta-Llama-3.1-8B-Instruct (Llama-3.1-8B), which runs locally through Ollama, generates a question, a short reference answer (the ground truth), and the supporting context (i.e., the source abstract itself). The identifier of its source paper was kept for every generated question. We verified that all the 80 stored contexts match their source abstracts exactly. % NK change 3: question-generation model named; workflow and settings referenced
The model was instructed to generate questions that can be answered directly using the information in the abstract. 
The evaluation workflow has been shown in Figure~\ref{fig:rag_eval} and its settings in Table~\ref{tab:rag_eval_setup}.

\paragraph{Question Review}
A human annotator (assisted by an AI tool; every decision was checked by the annotator) reviewed each question and its corresponding ground truth against the source abstract. Two simple rules were applied to determine the final list of questions used for evaluation:

\begin{itemize}
    \item Keep a question: A question is retained if it names its topic and its ground truth is supported by the source abstract. Flaws in a generated abstract, such as mixed topics, are not reason enough to remove a question because the abstract serves as the reference document. % NK fix: the rule is that the question names its topic, not that "its name" is supported by the abstract
    
    \item Remove a question: A question is removed if its ground truth is wrong, if it is too general to tell which of the 200 abstracts it refers to, or if it is wrongly worded.
\end{itemize}

% NK change 4: what the review did not use, and where the outcome is reported
The RAG answers were not considered during the review. The scores of the kept questions were recalculated from the original RAGAS results; no question was re-run. The outcome of the review is reported in Subsection~\ref{sec:rag_results}.

\subsection{System Evaluation}

So far, we have covered the system description, implementation details, and question generation. Now is the time to validate the effectiveness of our RAG-KBS. For this purpose, we utilize the Retrieval
Augmented Generation Assessment (RAGAS) framework \cite{es2024ragas}. Specifically, we use version 0.4.3 of RAGAS for evaluation. 
As shown in Figure~\ref{fig:RAG-Eval}, the end user sends a query via the FastAPI \cite{lathkar2023getting} web interface. 
The user query is embedded to vector representation. The remaining sub-steps continue in the background. Finally, the LLM replies with a grounded responses based on the most relevant documents. The RAGAS framework enables us to evaluate the quality of retrieved documents (i.e. the Abstracts in our case), and the grounded LLM response. In our application, rather than implementing RAGAS from scratch, we utilize the RAGAS open-source Python library \cite{ragas2024}. RAGAS offers reference-free evaluation metrics for RAG and LLM applications. However, of the four metrics used here, answer correctness and context recall compare the output with a reference answer. This is why the questions were generated together with their ground truths. These metrics fall into two categories: generation and retrieval metrics. All of these metrics are summarized in Table~\ref{tab:Ragas-Metrics}. % NK change 5: answer correctness and context recall are not reference-free

\begin{figure*}[ht]
\centering
\includegraphics[width=0.95\textwidth, height=0.45\textheight]{img/RAG-KBS.png}
\caption{The retrieval and generation part (online) from user query to grounded response. Note that the model names are subject to change. }
\label{fig:RAG-Eval}
\end{figure*}

% NK fix (typo): "subjective to change" -> "subject to change". NK answer to GN: I cannot see img/RAG-KBS.png; please check that the model names in it match Table rag_eval_setup (questions: Llama 3.1; RAGAS judge: Qwen2.5-14B; RAGAS embeddings: nomic-embed-text).

\begin{itemize}
    \item Generation Metrics: These metrics include faithfulness and answer correctness. Faithfulness refers to how factually accurate is the generated response with respect to the retrieved context. Whereas answer correctness refers to how well the generated response agrees with the ground truth. 
    % NK change 6: the evaluation used answer correctness (as in Table Ragas-Metrics), not answer relevance; faithfulness is measured against the retrieved context
    
    \item Retrieval Metrics: These metrics include context relevance (aka context precision) and context recall. Context relevance refers to the signal to noise ratio of retrieved context. On the other hand, context recall is meant to answer whether the LLM's retriever can retrieve all the necessary information (or ground-truth evidence) needed to answer a given question. % NK fix (typo): "weather" -> "whether". NK note (no change made): in RAGAS 0.4 "context relevance" and "context precision" are separate metrics; we calculated context precision only. Consider writing "context precision" here.
\end{itemize}

\begin{table}[H]
\caption{RAGAS evaluation metrics and retrieval measures.\label{tab:Ragas-Metrics}}
\begin{tabularx}{\textwidth}{lX}
\toprule
\textbf{Measure} & \textbf{What it checks} \\
\midrule
Faithfulness & Whether the claims in the answer are supported by the retrieved contexts \\
\midrule
Answer correctness & How well the answer agrees with the ground truth \\
\midrule
Context precision & Whether the relevant contexts are ranked near the top \\
\midrule
Context recall & Whether the retrieved contexts contain the information in the ground truth \\
\midrule
Hit@1 / Hit@20 & Whether the source abstract was ranked first / retrieved in the top 20 \\
\bottomrule
\end{tabularx}
\end{table}
% NK fix: Hit@1 and Hit@20 are not RAGAS metrics, so the caption now names both

%%\commentGN{WILL BE DELETED (or ELEVATED) AS QUESTION GENERATION IS MOVED IN 5.2}
% NK note (no change made): in the evaluated run, the questions were generated by Llama 3.1 and reviewed by a human annotator (Subsection Question Review); LLM-as-a-judge was not used to check the Q&A. If this paragraph is kept, please make it describe earlier attempts only.
%DEL%\commentRed{In addition, Ragas can automatically generate the questions and their corresponding ground truth answers (Q\&A dataset). For question synthesis prompt we utilized RAGAS's default SingleHopSpecificQuerySynthesizer template to construct factual benchmark question directly from sampled text without any manual additions. For answer generation prompt we used zero shot prompting because the model will not be hallucinated and follow a strict set of rules. We investigated this step based on the abstracts generated by the Llama-3.1-8b model and found it to a very time-consuming process. For instance, generating 10 questions took hours on a development laptop. As an alternative, we first prompted an LLM to generate the questions and answers (i.e., Q\&A). Then, we used the LLM-as-a-judge \cite{gu2026survey} approach to evaluate the quality of the resulting Q\&A. Another alternative for generating the Q\&A is to use Evidently AI platform \cite{evidentlyai}.}
%DEL%\commentRed{Having this Q\&A ready, we can evaluate our RAG-KBS using RAGAS framework, when it generates a response and provides a context for this response.}

% NK note (no change made): Figure~\ref{fig:rag_eval} (Methods) shows the pipeline that was actually executed; it could replace this figure.
\begin{figure*}[ht]
\centering
\includegraphics[width=0.95\textwidth, height=0.23\textheight]{img/rag_pipeline.pdf}
\caption{The RAGAS implementation pipeline.  } 
\label{fig:pipeline}
\end{figure*}

% NK change 7: the NK comment below is replaced by final text, which also addresses the Google Colab point

We evaluated the LLM grounded answers using four RAGAS metrics. In particular, we used Qwen2.5-14B (qwen2.5:14b) as the judge model and nomic-embed-text for embeddings. Both models were served by Ollama, and the evaluation was conducted in Google Colab on an NVIDIA Tesla T4 GPU (16~GB) memory, so no paid API was used. All metrics range from 0 to 1, where higher is better. Since every question carries the identifier of its source abstract, we also measured retrieval directly on the 20 stored contexts. For instance, Hit@1 is the share of questions whose source abstract was ranked first, and Hit@20 is the share whose source abstract appeared anywhere in the top 20.
Figure~\ref{fig:pipeline} summarizes the evaluation pipeline. After the configuration and output paths are set, the 200 generated abstracts are loaded from a CSV file and indexed in a Chroma vector database, one abstract per chunk. The local models are then initialised: Qwen2.5-14B as the judge and nomic-embed-text as the embedding model. Next, Llama-3.1-8B generates one question and a reference answer for each of 80 randomly selected abstracts. For each question, a dense vector search retrieves the 20 most similar abstracts, and Llama-3.1-8B writes an answer grounded in them. RAGAS then scores each answer on faithfulness, answer correctness, context precision and context recall. Finally, every question is validated against its source abstract, which leaves 52 valid questions, and the per-question scores, answers and test set are exported to CSV.

Following the approach provided by  \citet{aytar2024retrieval}, who focused on ``Context Relevance'', we can also focus on the same metric because the other metrics assess the quality of answers generated using the retrieved context. In our case, however, we consider the RAGAS metrics complementary to each other and report four metrics. 
%DEL%\ifshowchanges{\color{red}\textbf{[NK: deleted from here]}} \commentRed{(i) a pilot study that uses 50 questions, where the \textit{Qwen2.5-7B-Instruct} and \textit{Llama3.2:3b-Instruct} local models were used for generating the Q\&A dataset and grounded answer responses respectively; This configuration serves as a baseline for dense vector search with some context of lexical grounded text as well since pairing with hybrid retrieval strategy creates a lightweight, computationally efficient pipeline. (ii) the same number of questions,  where the \textit{Qwen2.5-14B-Instruct} and \textit{Llama3.1:8b-Instruct} local models were used to  generate the Q\&A dataset and grounded answer responses respectively;. Pairing with dense vector search this configuration's main purpose was to observe the performance semantic capacity matching. We used consistent hyperparameter setting and temperature = 0 for both the configurations to attain deterministic output. and (iii) a pilot study that uses 20 questions, where the \textit{GPT-4o mini} model was used for generating both Q\&A dataset and grounded answer responses. This enables us to perform a comparative system analysis as well. For the two settings, a sample of 50 questions was generated from 100 randomly selected abstracts on a consistent set of parameter setting.The configurations was selected because of the hardware restrictions and compatibility issues.} {\color{red}\textbf{[NK: to here]}}\fi
% NK change 8: executed configuration and hyperparameters (answer to GN). NK note for VP: please add the name of the answer-generation LLM, which is not in the files I have.
The reported evaluation was run once, in a single configuration (Table~\ref{tab:rag_eval_setup}): 80 questions generated by Llama~3.1 from 80 randomly selected abstracts; one abstract per chunk and the top 20 abstracts retrieved per question; RAGAS 0.4.3 with its default metric settings; the judge model Qwen2.5-14B with temperature 0 and at most 4,096 output tokens; and nomic-embed-text embeddings.


% NK answer to GN: the three items above are covered by the new Subsection "Evaluation Results" (sec:rag_results) at the end of this section: Table rag_scores gives the scores before and after the human review, and Tables rag_deleted, rag_kept_deleted and rag_examples describe its impact.

% NK change 9: pilot results removed, as marked by GN ("TO BE DELETED", "WILL BE CHANGED WITH RESULTS ANALYSIS"); they are replaced by Subsection sec:rag_results below

%DEL%\commentRed{Tables~\ref{tab:RagasL} and ~\ref{tab:RagasC} show the four different Ragas scores for our RAG-KBS using two samples that have 50 and 20 questions respectively. All of the reported metrics were obtained using an implementation similar to that in Figure~\ref{fig:Ragas}.}
% \begin{table}[H]
% \centering
% \caption{The reported Ragas scores (n=50). \commentRed{TO BE DELETED!!}}
% \label{tab:RagasL}
% \renewcommand{\arraystretch}{1.2}
% \begin{tabular}{lcc}
% \toprule
% \textbf{Metrics} & \multicolumn{2}{c}{\textbf{Score (\%)}} \\
% \cmidrule(lr){2-3}
% & \textbf{Qwen2.5:7b} & \textbf{Qwen2.5:14b} \\
% \midrule
% Faithfulness & 71.67 & 74.94\\
% Answer Relevance & 52.22  & 69.81\\
% Context Relevance/ Precision & 52.50  & 37.90\\
% Context Recall & 62.00 & 52.78\\
% \bottomrule
% \end{tabular}
% \end{table}

%DEL%\commentRed{The reported results in Table~\ref{tab:RagasL} relate to two models \textit{Qwen2.5:7b} and \textit{Qwen2.5:14b} respectively. For the \textit{Qwen2.5:7b} implementation, we used hybrid retriever. Thus, keyword matching was performed effectively hence we observe a solid context relevance/precision and context recall. Faithfulness was moderate,  showing that it was sometimes hallucinated, which costed it the metric score. LLAMA3.2:3b struggled on answer relevancy because of its smaller capabilities. Overall answer correctness is still solid because retrieved answer context helped the system to provide more factually correct answers. For the \textit{Qwen2.5:14b} implementation, we used a dense vector retriever, which added a huge amount of additional noise and irrelevant chunks. Consequently, we observed comparatively lower context relevance/precision and context recall. Since the model's capabilities were increased to \textit{LLAMA3.1:8b}, we noticed a marginal difference in faithfulness and answer relevance. It directly answered the question without any additional or unnecessary statements. Although the LLM was faithful and relevant, it was affected more by the context relevance and recall of the retrieval layer capping the overall answer correctness to just 43.25 \%.}

%DEL%% \commentRed{
% Whereas the reported results in Table~\ref{tab:RagasC} indicate strong grounding quality, with high faithfulness 95.83\% and context recall 90.00\%. These scores suggest that the generated answers are largely supported by the retrieved evidence. Moreover, the scores for context precision 85.35\% and answer relevance 80.70\% are slightly lower. This indicates that the system retrieves useful information. However, some responses contain minor irrelevances or excess context. Overall, the results confirm effective grounding, with the remaining limitations primarily in answer focus and conciseness. In this context, we want to mention that our RAG evaluation was conducted as a proof-of-concept using only 20 automatically generated questions and answers. Therefore, the RAGAS results should be interpreted as an initial validation rather than a comprehensive retrieval benchmark.
% }

% \begin{table}[H]
% \centering
% \caption{The reported Ragas scores (n=20).}
% \label{tab:RagasC}
% \renewcommand{\arraystretch}{1.2}
% \begin{tabular*}{0.55\textwidth}{@{\extracolsep{\fill}}lr}
% \toprule
% \textbf{Metric} & \textbf{Score (\%)} \\
% \cmidrule(lr){2-2}
% & \textbf{GPT-4o mini} \\
% \midrule
% Faithfulness & 95.83\\
% \midrule
% Answer Relevance & 80.70\\
% \midrule
% Context Relevance/ Precision & 85.35\\
% \midrule
% Context Recall & 90.00\\
% \bottomrule
% \end{tabular*}
% \end{table}

% ===================== NK: new subsection with the RAG evaluation results =====================
% NK change 10: results of the RAG evaluation (reviewer comments 2, 7, 8, 9, 11, 12, 14, 15); fills GN's "before / human factor / after" placeholders

\subsection{Evaluation Results}
\label{sec:rag_results}

\paragraph{Question validation}
Of the 80 generated questions, 52 (65.0\%) were kept and 28 (35.0\%) were deleted (Table~\ref{tab:rag_deleted}). Most deletions, 25 of 28, were questions that do not name their topic, such as ``What is the primary goal of the research presented in this paper?''; two questions had a wrong reference answer and one was wrongly worded. RAGAS returned scores for 78 of the 80 questions: questions 55 and 64 have the same wording as question 4, and the evaluation script identified completed questions by their text, so it skipped them. Both are among the deleted questions, so all 52 kept questions have scores. Table~\ref{tab:rag_examples} gives examples of kept and deleted questions, and Appendix~\ref{app:sub_questions} lists a larger sample.

\begin{table}[H]
\centering
\caption{List of deleted evaluation questions by reason.}
\label{tab:rag_deleted}
\footnotesize
\begin{tabularx}{\textwidth}{p{0.28\textwidth} c X}
\toprule
\textbf{Reason} & \textbf{Count} & \textbf{Questions} \\
\midrule
Question too general & 25 & Q1, Q4, Q5, Q7, Q17, Q18, Q22, Q27, Q30, Q33, Q34, Q40, Q42, Q43, Q44, Q51, Q53, Q55, Q57, Q60, Q63, Q64, Q66, Q69, Q77 \\
\midrule
Reference answer wrong & 2 & Q32, Q76 \\
\midrule
Question wrongly worded & 1 & Q21 \\
\midrule
Total & 28 & \\
\bottomrule
\end{tabularx}
\end{table}


\paragraph{RAGAS scores before and after validation}
All four metrics were higher on the 52 valid questions than on all 78 scored questions (Table~\ref{tab:rag_scores}, Figure~\ref{fig:rag_results}a). Context precision and answer correctness gained the most; context recall changed least. After validation, 43 of 52 questions received the maximum faithfulness score and 44 of 52 the maximum context precision score.


\begin{table}[H]
\centering
\caption{The RAGAS scores before (including all scored questions) and after (including valid questions only) question validation: mean $\pm$ SD with 95\% bootstrap confidence interval; $n$ is the number of questions.}
\label{tab:rag_scores}
\footnotesize
\setlength{\tabcolsep}{3pt}
\begin{tabularx}{\textwidth}{X c c c c c c c}
\toprule
\textbf{Metric} & $\boldsymbol{n}$ & \textbf{Before} & \textbf{95\% CI} & $\boldsymbol{n}$ & \textbf{After} & \textbf{95\% CI} & \textbf{Change} \\
\midrule
Faithfulness & 77 & 0.777 $\pm$ 0.395 & 0.686--0.862 & 52 & 0.892 $\pm$ 0.270 & 0.814--0.959 & +0.116 \\
Answer correctness & 78 & 0.569 $\pm$ 0.283 & 0.507--0.631 & 52 & 0.689 $\pm$ 0.225 & 0.628--0.748 & +0.121 \\
Context precision & 78 & 0.778 $\pm$ 0.378 & 0.693--0.860 & 52 & 0.935 $\pm$ 0.186 & 0.880--0.980 & +0.157 \\
Context recall & 78 & 0.720 $\pm$ 0.438 & 0.623--0.811 & 52 & 0.753 $\pm$ 0.414 & 0.637--0.860 & +0.033 \\
\bottomrule
\end{tabularx}
\noindent{\footnotesize{*Faithfulness before validation uses 77 questions because question 27 received no answer, so there were no claims to check.}}
\end{table}

\begin{figure}[H]
\centering
\includegraphics[width=\textwidth]{img/rag_evaluation_results.pdf}
\caption{Results of the RAG evaluation. (\textbf{a}) Mean RAGAS scores on all 78 scored questions and on the 52 valid questions; error bars show 95\% bootstrap confidence intervals. (\textbf{b}) Share of questions whose source abstract was ranked first (Hit@1) or retrieved among the top 20 (Hit@20), for all questions, valid questions and the deleted general questions.}
\label{fig:rag_results}
\end{figure}

\paragraph{Valid versus deleted questions}
The deleted questions scored significantly lower than the valid ones on faithfulness, answer correctness and context precision (Table~\ref{tab:rag_kept_deleted}), which supports H2. Context recall did not differ between the two groups.

\begin{table}[H]
\centering
\caption{RAGAS scores of valid and deleted questions: mean $\pm$ SD ($n$); two-sided Mann--Whitney U test.}
\label{tab:rag_kept_deleted}
\footnotesize
\begin{tabularx}{\textwidth}{X c c c}
\toprule
\textbf{Metric} & \textbf{Valid} & \textbf{Deleted} & \textbf{$p$-value} \\
\midrule
Faithfulness & 0.892 $\pm$ 0.270 (52) & 0.536 $\pm$ 0.499 (25) & 0.001* \\
Answer correctness & 0.689 $\pm$ 0.225 (52) & 0.327 $\pm$ 0.228 (26) & $<$0.001* \\
Context precision & 0.935 $\pm$ 0.186 (52) & 0.463 $\pm$ 0.463 (26) & $<$0.001* \\
Context recall & 0.753 $\pm$ 0.414 (52) & 0.654 $\pm$ 0.485 (26) & 0.46 \\
\bottomrule
\end{tabularx}
\noindent{\footnotesize  $^*$Statistically significant ($p < 0.05$).}
\end{table}

\paragraph{Retrieval of the source abstract}
Across all 80 questions, the source abstract was ranked first for 60 (75.0\%) and retrieved among the top 20 for 64 (80.0\%). For the 52 valid questions, these rates were 94.2\% and 98.1\%; only one valid question (Q52) did not retrieve its source abstract. For the 25 questions that were too general, the source abstract was ranked first for 8 (32.0\%) and retrieved for 10 (40.0\%) (Figure~\ref{fig:rag_results}b). The difference between valid and general questions is significant for both Hit@1 and Hit@20 (Fisher's exact test, $p < 0.001$), which supports H1.

\paragraph{Context recall}
Context recall was 0 for 11 of the 52 valid questions, although for 10 of them the source abstract had been ranked first. For example, Q65 (Table~\ref{tab:rag_examples}) was answered correctly from its top-ranked source abstract (answer correctness 0.952) but received a context recall of 0. In our evaluation, context recall therefore did not reflect retrieval quality, which explains why it changed little after validation; Hit@$k$, calculated from the known source identifiers, gave a clearer picture of retrieval.


\paragraph{Computational cost}
Measured from the time stamps saved by the evaluation script, scoring one question with the four RAGAS metrics took a median of 3.3 minutes on the Tesla T4 GPU (mean 3.5, range 2.2--6.7 minutes; 72 consecutive questions, excluding 5 interruptions of the run longer than 60 minutes), or an estimated 4.6 hours for the 78 scored questions. All models ran locally, so the evaluation incurred no API fees.

\begin{table}[H]
\centering
\caption{Examples of valid and deleted evaluation questions. Long answers are shortened (marked using \ldots).}
\label{tab:rag_examples}
\scriptsize
\begin{tabularx}{\textwidth}{p{0.6cm} X p{2.2cm} X p{2.4cm}}
\toprule
\textbf{Q} & \textbf{Question} & \textbf{Reference answer} & \textbf{RAG answer} & \textbf{Decision} \\
\midrule
Q65 & What percentage of proteins exhibit complex knotting topologies? & approximately 10\% & Approximately 10\%. (Paper ID: 704.0191) &Valid; answer correct, but context recall 0 \\
\midrule
Q11 & What type of distribution do the time intervals between flag flaps follow? & An exponential distribution & The provided abstracts do not contain enough information \ldots{} according to Paper ID: 704.0484, \ldots{} ``the time intervals between flaps follow an exponential distribution''. & Valid; answer correct but hedged (faithfulness 0.500) \\
\midrule
Q4 & What is the primary goal of the research presented in this paper? & Extracting text from historical handwritten documents & The provided abstracts do not contain enough information to determine which paper you are referring to \ldots & Deleted: too general \\
\midrule
Q32 & What is the primary role of electromagnetic field enhancement in the PFE phenomenon? & Mie resonance and Bragg diffraction & The primary role of electromagnetic field enhancement due to Mie resonance and Bragg diffraction plays a crucial role in the photonic flame effect (PFE) phenomenon \ldots & Deleted: reference answer gives the cause, not the role \\
\midrule
Q21 & Does the RAG system change space topology with respect to Maxwell's equations? & Yes & Yes. The provided abstracts do not contain enough information about the ``RAG system'', but \ldots & Deleted: wrongly worded \\
\bottomrule
\end{tabularx}
\end{table}

%%%%%%%%%%%%%%%%%%%%%%%%%%%%%%%%%%%%%%%%%%
\section{Discussion}

%DEL%\commentGN{Research Summary --> Interpretation of findings -- comparison with literature (optional) --> implications of work --> limitations (Point Wise) --> Future Work}

% NK change 1: research summary written for the executed evaluation (80 generated / 52 valid questions; answer correctness instead of answer relevancy). The VP paragraph below described the 50-question pilot, whose tables were deleted from Section KB as GN marked.

\textbf{Research summary.} This study generated abstracts for 200 arXiv papers using three small, locally deployed LLMs, indexed the generated abstracts in a RAG-based KB, and evaluated the KB with LLM-generated questions that were validated by a human annotator. Llama-3.1-8B produced the shortest abstracts and the highest scores on all automatic metrics (Table~\ref{tab:allresults}). For the KB, 52 of the 80 generated questions were valid. On these questions, the KB reached a mean faithfulness of 0.892, answer correctness of 0.689, context precision of 0.935 and context recall of 0.753, and it ranked the source abstract first for 94.2\% of the questions (Section~\ref{sec:rag_results}).

% NK change 2: interpretation of the executed results. The four VP paragraphs below compared the pilot configurations (Qwen2.5 7B vs 14B, hybrid vs dense retrieval); those results are no longer in the paper, so the paragraphs are deleted. The hybrid-retrieval idea is kept as future work.
% \ifshowchanges{\color{red}\textbf{[NK: deleted from here]}}
% \commentVP{When comparing faithfulness, it suggests \commentRed{that both have the grounded quality response and the marginal gain by 14b \& 8b configuration parameter model is because of its higher ability for parametric synthesis. Similarly, for answer relevancy it highlighted the configuration gap as 14b \& 8b configuration model is able to answer more on topic and has the capacity filter out irrelevant chunks from abstract.}}

% \commentVP{For context relevancy/precision and context recall,\commentRed{ since we use a hybrid approach it outperforms dense vector search. This only suggests one thing that BM25 was able to retrieve more ground-truth relevant lexical text and exact matching words to output a higher result even though dense vector search performs better result factually, it is not necessary that it will output answer with high critical vocabulary constraints and precision.}}

% \commentVP{From these we can observe that \commentRed{ dense retrieval + higher model performs better on generation, though hybrid retrieval + smaller model dominates on strictly sticking to the point and context.}}

% \commentVP{The findings and results of this study imply that \commentRed{regardless we doubled the size of the model a hybrid retrieval is still more robust. Hence, integrating with hybrid retrieval with a higher model with improve it even further.}}
% {\color{red}\textbf{[NK: to here]}}\fi

\textbf{Interpretation of findings.} The main finding of the RAG evaluation is that the quality of the generated questions strongly affects the measured quality of the system. More than one third of the generated questions (28 of 80) were invalid, and 25 of them were too general to identify their source abstract, e.g., ``What is the primary goal of the research presented in this paper?''. Such a question is clear within the abstract it was generated from, but ambiguous in a corpus of 200 abstracts. The retrieval results show this directly: the source abstract was retrieved for 98.1\% of the valid questions, but for only 40.0\% of the general ones (H1). When the source abstract is missing, the answer cannot be supported by the context, so the deleted questions also scored lower on faithfulness, answer correctness and context precision (H2). In other words, much of the gap between the scores before and after validation reflects flaws in the questions rather than in the KB. Context recall was the exception: it changed little after validation and did not differ between valid and deleted questions, because the judge often scored it 0 even when the source abstract was ranked first (10 of the 52 valid questions). In our setting, Hit@$k$ was therefore a more reliable measure of retrieval than context recall. A wrong reference answer also penalizes a correct RAG answer: for Q32, the answer agreed with the source abstract, but it received an answer correctness of only 0.385 because the reference answer was wrong. Among the valid questions, answer correctness (0.689) remained the lowest score although retrieval was nearly perfect. This suggests that the remaining errors arise mainly in answer generation, for example hedged answers such as Q11 in Table~\ref{tab:rag_examples}. However, note that answer correctness also includes a semantic-similarity component, and no question reached the maximum score of 1.

% NK fixes (grammar): "main point" -> "main points"; "carried" -> "carried out"; "pipeline which are" -> "pipelines, which are"
One of the main points of this study is that all of the experiments were carried out locally since standard RAG pipelines are cloud-based and expose data to third parties. However, our approach highly aligns with \citet{atkinson2026deterministic}'s execution, since they strictly follow the local constraints and zero data egress, which is also required for high stakes institutional AI deployments similar to this one.
% NK change 3: implications of the work
The RAG evaluation was also run with local models only: scoring one question with the four RAGAS metrics took a median of 3.3 minutes on a single Tesla T4 GPU, about 4.6 hours for all scored questions, with no API fees. Local evaluation is therefore feasible, at the cost of longer run times.

\textbf{Implications.} Firstly, LLM-generated evaluation questions should be validated against their source documents before they are used, and RAGAS scores should be reported on the validated set; otherwise the scores mix the quality of the system with the quality of the test set. Additionally, a question generator that sees one document at a time should be instructed to name the topic of the document in every question, because general questions were the most common error. 
Also, because each generated question keeps the identifier of its source document, retrieval can be measured directly with Hit@$k$, which needs no judge model and should be reported alongside RAGAS. Hit@$k$ also helps to find invalid questions: a valid question whose source document is not retrieved deserves a closer look.
Lastly, generated corpora should be cleaned of introductory text before indexing, and evaluation scripts should track questions by a unique identifier rather than by their text, since identical wording caused two questions to be skipped.

% NK change 4: limitation points replaced, as GN marked. The old points described the pilot (50 questions, SingleHopSpecificQuerySynthesizer, same model for generation and judging) and a CS/ML-only corpus, which no longer match the paper: the corpus covers several fields (arXiv IDs 0704.0002-0704.1627; reviewer comment 4), and questions were generated by Llama 3.1 and judged by Qwen2.5-14B.
\paragraph{Limitations}

\begin{itemize}
    \item \textit{Synthetic data.} The abstracts, the questions, the reference answers and the metric judgments were all produced by LLMs; the KB was not evaluated against an established external benchmark.

    \item \textit{Single configuration.} One retrieval configuration (one abstract per chunk, top-$k = 20$), one answer-generation model and one judge model were tested. 
    %Note:OH%\commentRed{Rooms for improvement?}

    \item \textit{Corpus.} The KB contains 200 arXiv abstracts from a single month (April 2007). The results may not generalize to larger corpora, other periods or other types of documents.

    \item \textit{Number of questions.} The RAG evaluation is based on 80 generated questions, of which 52 are valid. The confidence intervals in Table~\ref{tab:rag_scores} are therefore wide, e.g., 0.637--0.860 for context recall.

    \item \textit{Question type.} Each question was generated from a single abstract, so all questions are single-hop factual questions. Questions that need several documents or comparisons were not evaluated.

        \item \textit{Question review.} The questions were reviewed by one annotator, assisted by an AI tool, so inter-annotator agreement could not be measured; a second independent reviewer would make the review more robust.

    \item \textit{Judge model.} The RAGAS scores come from a single LLM judge (Qwen2.5-14B) with the default RAGAS settings; their agreement with human judgments was not measured. The evaluation was run once, so run-to-run variation was not measured either.

    \item \textit{Missing scores.} Two deleted questions (Q55 and Q64) were not scored because they have the same wording as Q4 and the evaluation script identified completed questions by their text. All valid questions were scored, so the results after validation are not affected.

\end{itemize}

% NK change 5: future work, following GN's structure
\textbf{Future work.}  We plan to (i) extend the evaluation to a larger corpus and more questions; (ii) prompt the question generator to name the topic of each question and to produce multi-hop questions; (iii) have the questions reviewed by several annotators and measure their agreement; (iv) compare the judge scores with human judgments and with other judge models; and (v) compare retrieval configurations, such as dense, hybrid (BM25 and dense) and reranked retrieval, and answer-generation models of different sizes on the same validated question set.

%%%%%%%%%%%%%%%%%%%%%%%%%%%%%%%%%%%%%%%%%%
\section{Conclusion}
\label{sec:conclusion}
%\commentRed{Following IMRaD structure, the Conclusions can be merged with the Discussion!}
In this study, we implemented a RAG framework for institutional knowledge management using locally served LLMs and evaluated it using the RAGAS evaluation framework. 
This research aimed at automating the generation of abstracts to enable semantic search within large, unstructured document repositories while maintaining data privacy. 
We utilized the LM Studio platform to evaluate three locally deployed: Mixtral-8$\times$7B, Qwen3-VL-8B, and Llama-3.1-8B on a subset of 200 technical papers from the arXiv dataset. Our methodology involved  three steps: (i) a comparative evaluation of model performance;  (ii) the development of a KB system employing Ollama, Chroma VDB and an abstract-first retrieval approach , and (iii) the evaluation of the KB with RAGAS on LLM-generated questions that were validated against their source abstracts. % NK change 1: third step for the RAG evaluation (reviewer comments 3, 5); "two-steps" -> "three steps" (also removes the wrong hyphen). NK note for VP: please confirm Chroma VDB (see the note in Section KB).

Our experimental findings showed that Llama-3.1-8B consistently outperformed the other small LLMs across standard lexical and semantic metrics, including BLEU, ROUGE, and BERTScore. Furthermore, the ablation studies revealed that decoding parameters, such as temperature and sampling thresholds, just have a limited impact on the quality of the generated abstracts. However, fully structured, instruction-based zero-shot prompts are essential for generating coherent and concise summaries. Overall, these results indicate that small, quantized models deployed on local hardware offer a robust, privacy-preserving solution for managing institutional technical documentation. With  regard to ITE, following the theoretical and experimental results, the near-term future work will involve transitioning the RAG-KBS from a demonstration phase to a functional implementation on the ITE infrastructure. This implementation will be specifically tailored to automate knowledge management for secure, long-term research projects. % NK fixes (grammar): missing comma after "thresholds"; missing article "a"; "With regards to" -> "With regard to"; sentence fragment ". The" -> ", the"

% NK change 2: conclusions of the RAG evaluation (reviewer comments 2, 8, 14, 15)
The evaluation of the RAG-based KB showed that LLM-generated evaluation questions must be validated before use. Of the 80 questions generated by Llama-3.1-8B, 28 (35.0\%) were invalid, mostly because they were too general to identify their source abstract. The general questions retrieved their source abstract in only 40.0\% of cases, against 98.1\% for the valid questions, and the invalid questions scored significantly lower on faithfulness, answer correctness and context precision. Once the invalid questions were removed, all four RAGAS scores of the RAG-based KB improved, most clearly context precision and answer correctness, while context recall changed little. The whole evaluation ran on local models on a single Tesla T4 GPU at a few minutes per question and with no API fees. We therefore recommend that RAG systems evaluated with synthetic questions report their scores on a validated question set.

%In analogy to ITE Small locally deployed LLMs can support text summarization and abstracts generation for institutional archives when prompts and decoding settings are carefully controlled.
%This section is not mandatory, but can be added to the manuscript if the discussion is unusually long or complex.
%%%%%%%%%%%%%%%%%%%%%%%%%%%%%%%%%%%%%%%%%%
\section{Further Details}
%%%%%%%%%%%%%%%%%%%%%%%%%%%%%%%%%%%%%%%%%%
\authorcontributions{Conceptualization, O.H.; methodology, O.H.; Literature review and related work analysis T.D.; software and implementation, O.H., N.Z., and V.P.; validation, O.H., Z.K. and V.P.; formal analysis, O.H. and N.Z.; investigation, O.H.; resources, G.N. and O.H.; data curation, O.H.; original draft preparation, O.H.; writing---review and editing, O.H., T.D., N.Z., and V.P.; visualization, O.H. and N.Z.; supervision, O.H. and G.N.; administration, G.N.; funding acquisition, G.N.; all authors have read and agreed to the submitted version of the manuscript.}

\funding{This research received no external funding.}

\dataavailability{The data presented in this study is a subset of the arXiv dataset. It is available upon request from the
corresponding author. It will also be made available in the published version. The Code, released resources, and the artifacts are available at the GitHub repository.\\ \url{https://github.com/ohamed/lmstudio}.
}

% NK note (no change made): this section is empty; please fill it or remove the heading before submission.
\acknowledgments{
During the preparation of this manuscript, the authors used some generative artificial intelligence (GenAI) tools for limited and non-substantive purposes. ChatGPT were used to assist English language editing and grammar refinement in Section~\ref{sec:related_work}. DeepL Pro was also used elsewhere to refine and sometimes restructure the English text, and to translate reference \cite{bektacs2025academic} from Turkish to English. Napkin AI was used to create the data preprocessing workflow diagram shown in Figure~\ref{fig:Prep-prep}. Furthermore, the free version of ChatGPT Images (or GPT-Image) was used to create Figure~\ref{fig:articles}, whereas Nano Banana was used to create the diagrams in Figure~\ref{fig:system_architecture} and Figure~\ref{fig:RAG}. In addition, ChatGPT and Perplexity AI were consulted to assist in outlining aspects of the experimental design related to the ablation studies.
%DeepL was used to help restructure the English in different sections.

The authors also confirm that all scientific content, methodology, analysis, and conclusions are their own and take full responsibility for the integrity of the work. 
%DEL:% The authors would like to thank \textbf{Mr. Vibhav Patil}, our student and research assistant for taking part in implementing the RAGAS results, particularly those  in Table 8.
}


\conflictsofinterest{The authors declare no conflicts of interest. 
%Authors must identify and declare any personal circumstances or interest that may be perceived as inappropriately influencing the representation or interpretation of reported research results. Any role of the funders in the design of the study; in the collection, analyses or interpretation of data; in the writing of the manuscript; or in the decision to publish the results must be declared in this section. If there is no role, please state ``The funders had no role in the design of the study; in the collection, analyses, or interpretation of data; in the writing of the manuscript; or in the decision to publish the results''.
} 

%%%%%%%%%%%%%%%%%%%%%%%%%%%%%%%%%%%%%%%%%%
% \abbreviations{Abbreviations}{
% The following abbreviations are used in this manuscript:
% \\
% \noindent 
% \begin{tabular}{@{}ll}
% MDPI & Multidisciplinary Digital Publishing Institute\\
% DOAJ & Directory of open access journals\\
% TLA & Three letter acronym\\
% LD & Linear dichroism \\
% LLMs & Large Language Models
% \end{tabular}
% }

%%%%%%%%%%%%%%%%%%%%%%%%%%%%%%%%%%%%%%%%%%
%% Optional
\appendixtitles{no} 
% Leave argument "no" if all appendix headings stay EMPTY (then no dot is printed after "Appendix A"). If the appendix sections contain a heading then change the argument to "yes".

%%%%%%%%%%%%%%::App§A§
\appendixstart
\appendix
\section[\appendixname~\thesection]{}
\subsection[\appendixname~\thesubsection]{Software and Hardware}
\label{app:sub_sw_hw}
This appendix briefly describes the necessary combination of software and hardware required to run our experiments  (i.e., part of experimental setup). It also provides configuration details.

%\subsubsection{Software and Hardware}
\paragraph{Software} The necessary software for starting this experimental setup falls into three categories.

\begin{itemize}
\item	Software Packages: The LM Studio \textit{Version 0.3.31} or higher must be installed on hardware with sufficient computational power.

\item	Local LLMs: We utilized three local LLMs that belong to different families. Their corresponding quantization file names are: \textbf{mixtral\_instruct\_7b.Q8.gguf}, \textbf{qwen3-vl-8b-instruct-abliterated-q4-k-m.gguf}, and \textbf{Meta-llama-3.1-8b-instruct-Q4\_K\_S.gguf}. Additionally, for the LM Studio GGUF runtime selections, we used \textbf{CUDA llama.cpp (Windows) v1.52.1}. For more information about these models, visit their web pages on Hugging Face (see Appendix~\ref{app:sub_hf}). % NK fix (typo): LMMs -> LLMs. NK note (no change made): the paper calls the first model Mixtral-8$\times$7B, but the file name (mixtral_instruct_7b) and the Hugging Face page (LeroyDyer/Mixtral_Instruct_7b) suggest a 7B model. Authors of this part: please check which name is correct.

\item	Python Libraries: Creating a virtual environment using the Python built-in \textbf{venv} module is a prerequisite for this part. Being done with that, the needed Python libraries can be installed using the Python PIP, which is a package manager for installing the Python libraries. % NK note (no change made): "Being done with that" could read "Once this is done".
% The following libraries were installed in the created virtual environment:
The list of Python dependency libraries is provided in the \textbf{requirements.txt} file. These libraries should be installed in the created virtual environment.
\end{itemize}

%\subsubsection{Hardware}
\paragraph{Hardware}
The abstract generation experiments were carried out on a Windows machine (MS Windows 11 installed), which is equipped with an NVIDIA RTX PRO 3000 Blackwell Generation laptop GPU. The RAG evaluation was carried out on different hardware, described below.

\paragraph{Local API Configuration}
All the three local models were accessed via the REST API call endpoint on the local machine, specifically \url{http://localhost:1234/v1/}.

\paragraph{LM Studio Configuration} The LM Studio configuration presets are JSON files. They are located in the GitHub repository within the \textit{config-presets} folder.

% NK change 1: software and hardware of the RAG evaluation (reviewer comments 3, 6; GN: "The Google Colab point should be addressed")
\paragraph{RAG Evaluation Setup.} The RAG evaluation (Section~\ref{sec:KB}) did not use LM Studio. The RAGAS scoring was run in Google Colab on an NVIDIA Tesla T4 GPU (15,360 MiB). 
The judge model Qwen2.5-14B (qwen2.5:14b) and the embedding model nomic-embed-text were served by Ollama and called through its OpenAI-compatible API alongside  the Python libraries ragas 0.4.3 and OpenAI 1.109.1. 
The judge used temperature 0 and at most 4,096 output tokens. 
The four metrics (Faithfulness, AnswerCorrectness, ContextPrecision and ContextRecall) were imported from \texttt{ragas.metrics.collections} and used with their default settings. The evaluation questions were generated with Llama~3.1 (llama3.1), also run through Ollama. Table~\ref{tab:rag_eval_setup} summarizes these settings.

\subsection[\appendixname~\thesubsection]{Hugging Face}
\label{app:sub_hf}
%\paragraph{\textbf{Hugging Face}}
Further details about the three local models are available on the Hugging Face platform.

\begin{itemize}
\item Visit \textbf{Mixtral\_Instruct\_7b} web pages at {\scriptsize \url{https://huggingface.co/LeroyDyer/Mixtral_Instruct_7b}}.
\item Visit \textbf{Qwen3-vl-8b-instruct} web pages at {\scriptsize \url{https://huggingface.co/wangkanai/qwen3-vl-8b-instruct}}.
\item Visit \textbf{Meta-Llama-3.1-8B-Instruct-GGUF} web pages at {\scriptsize \url{ https://huggingface.co/tensorblock/Meta-Llama-3.1-8B-Instruct-GGUF}}.
\end{itemize}

%%%%%%%%%%%%%%::App§B§
% \appendixstart
%\appendix
\section[\appendixname~\thesection]{}
\subsection[\appendixname~\thesubsection]{Prompt Designs}
\label{app:sub_prompt_des}

This appendix is meant to provide further details that are crucial to understanding, and reproducing the research shown in this paper.

\begin{table}[H]
%\centering
\small
%\footnotesize
\caption{The prompts used for the ablation  study in LLM.
\label{tab:prompts-design}}
\resizebox{\textwidth}{!}{
%\newcolumntype{C}{>{\centering\arraybackslash}X}
\begin{tabularx}{\textwidth}{llX}
\toprule
\textbf{ID}	& \textbf{Prompt Type}	& \textbf{Prompt Text}\\
\midrule
P1		& Zero-shot	& Write a formal academic abstract based on the following paper content:

\{\}\\
\midrule
P2		& Zero-shot	& Write a formal academic abstract in one continuous paragraph based on the following paper content:

\{\}

Word count:\\
\midrule
P3		& Zero-shot	& Write a formal academic abstract of 120-130 words in one continuous paragraph based on the following paper content:

\{\}

Word count:\\
\midrule
P4		& Zero-shot	& Write a formal academic abstract of 120-130 words in one continuous paragraph. Cover the background, methods, main results, and conclusions. Base it on the following paper content:

\{\}

Word count:\\
\midrule
P5		& Instruction-based (Zero-shot)	& Write ONLY a new formal academic abstract of exactly 120--130 words in one continuous paragraph.
Strictly cover: (1) background or motivation, (2) methods or approach,
(3) main results, (4) conclusions or implications.
Do NOT include feedback, suggestions, critiques, numbered lists, or references to ``your abstract''.
Base the abstract solely on the provided paper content.\\
\bottomrule
\end{tabularx}
}
\end{table}

%%%%%%%%%%%%%%::App§C§
% \appendixstart
%\appendix
\section[\appendixname~\thesection]{}
% \subsection[\appendixname~\thesubsection]{Generated Questions}
% \label{app:sub_gen_ques}

% NK change 2: sample of the generated evaluation questions (answer to the request above; reviewer comments 2, 11)

\subsection[\appendixname~\thesubsection]{Sample of the Generated Evaluation Questions}
\label{app:sub_questions}

Table~\ref{tab:app_questions} lists a sample of the 80 questions generated by Llama~3.1 for the RAG evaluation (Section~\ref{sec:KB}). To avoid selecting favourable examples, the sample contains the first 20 questions in the order in which they were generated (Q1--Q20), together with the three questions deleted for reasons other than being too general (Q21, Q32 and Q76). Each question is shown with its reference answer, the rank of its source abstract among the 20 retrieved abstracts, and the outcome of the question review. Questions and reference answers are given as generated.

\begin{table}[H]
\caption{Sample of the generated evaluation questions. \textbf{Rank}: position of the source abstract among the 20 retrieved abstracts (--: not retrieved).}
\label{tab:app_questions}
\scriptsize
\renewcommand{\arraystretch}{1.15}
\begin{tabularx}{\textwidth}{@{}p{0.55cm} p{1.25cm} >{\raggedright\arraybackslash}X >{\raggedright\arraybackslash}p{3.3cm} c >{\raggedright\arraybackslash}p{1.9cm}@{}}
\toprule
\textbf{Q} & \textbf{arXiv ID} & \textbf{Question} & \textbf{Reference answer} & \textbf{Rank} & \textbf{Decision} \\
\midrule
Q1 & 0704.0712 & Does the study use a simplified model potential? & yes & -- & Deleted: too general \\
Q2 & 0704.0071 & According to the abstract, what percentage of typological differences between languages indicates they are unlikely to be related? & 0.6 & 1 & Valid \\
Q3 & 0704.0157 & What type of systems does the proposed approach provide equations of state (EOS) for? & single-component and multi-component hard-sphere fluids & 1 & Valid \\
Q4 & 0704.1267 & What is the primary goal of the research presented in this paper? & Extracting text from historical handwritten documents & -- & Deleted: too general \\
Q5 & 0704.0967 & What is the specific problem that the researchers are trying to address? & The researchers are jointly optimizing DPC power allocation and multihop/multipath routing in MIMO-based wireless mesh networks. & -- & Deleted: too general \\
\midrule
Q6 & 0704.0860 & What is the duration of unavailability estimated by the proposed method for servers? & 12 hours per year & 1 & Valid \\
Q7 & 0704.0499 & Does the proposed heuristic algorithm run in polynomial time? & Yes & -- & Deleted: too general \\
Q8 & 0704.1334 & What is the primary purpose of developing analog electronics for serial readout of silicon strip sensors? & High precision vertex detectors around high energy colliders & 1 & Valid \\
Q9 & 0704.1373 & Does Zebu's added safety and robustness significantly impact performance? & No & 1 & Valid \\
Q10 & 0704.0301 & What type of techniques are used to develop the framework for single- and multi-place differentially algebraic functions? & A combination of analytic and algebraic techniques & 1 & Valid \\
\midrule
Q11 & 0704.0484 & What type of distribution do the time intervals between flag flaps follow? & An exponential distribution & 1 & Valid \\
Q12 & 0704.1513 & What is the sufficient value of gain required for complete compensation of total loss in low-loss silver film and low-index dielectric? & approximately 420 cm$^{-1}$ & 1 & Valid \\
Q13 & 0704.1314 & What type of analysis was used to investigate acoustic phenomena in microfluidic chambers? & Full-image micro-PIV analysis & 1 & Valid \\
Q14 & 0704.0575 & What type of potentials satisfy the conditions for bounded and closed orbits? & The Newtonian potential and the isotropic harmonic oscillator potential. & 1 & Valid \\
Q15 & 0704.1542 & What theoretical framework does the research employ to calculate single-particle Green's functions? & The propagator theory framework & 1 & Valid \\
\midrule
Q16 & 0704.1411 & What type of codes are used as an alternative approach to traditional TCQ systems? & Maximum-Hamming-distance convolutional codes & 1 & Valid \\
Q17 & 0704.0363 & What is the main concept that this paper aims to demonstrate? & The invariance of measured length of a ruler in motion & -- & Deleted: too general \\
Q18 & 0704.1196 & What is the worst-case time complexity of the proposed algorithm? & $O((d^2)n)$ & 1 & Deleted: too general \\
Q19 & 0704.0598 & What type of network does the ENG model evolve into? & A scale-free graph & 2 & Valid \\
Q20 & 0704.0492 & What type of attack proposed by Liu and Zhang is considered a pseudo-attack? & Their private key attack & 1 & Valid \\
\midrule
Q21 & 0704.0914 & Does the RAG system change space topology with respect to Maxwell's equations? & Yes & 1 & Deleted: wrongly worded \\
Q32 & 0704.0088 & What is the primary role of electromagnetic field enhancement in the PFE phenomenon? & Mie resonance and Bragg diffraction & 1 & Deleted: reference answer wrong \\
Q76 & 0704.0856 & Do radiative transitions with $\Delta l = -1$ dominate numerically everywhere? & No, except for small values of eccentricity $\varepsilon$ close to unity. & 1 & Deleted: reference answer wrong \\
\bottomrule
\end{tabularx}
\end{table}

% % All appendix sections must be cited in the main text. In the appendices, Figures, Tables, etc. should be labeled, starting with ``A''---e.g., Figure A1, Figure A2, etc.

%%%%%%%%%%%%%%%%%%%%%%%%%%%%%%%%%%%%%%%%%%
%\isPreprints{}{% This command is only used for ``preprints''.
\begin{adjustwidth}{-\extralength}{0cm}
%} % If the paper is ``preprints'', please uncomment this parenthesis.
%\printendnotes[custom] % Un-comment to print a list of endnotes

\reftitle{References}

% Please provide either the correct journal abbreviation (e.g. according to the “List of Title Word Abbreviations” http://www.issn.org/services/online-services/access-to-the-ltwa/) or the full name of the journal.
% Citations and References in Supplementary files are permitted provided that they also appear in the reference list here. 

%=====================================
% References, variant A: external bibliography
%=====================================
\bibliography{references}

%=====================================
% References, variant B: internal bibliography the "ACS format".
% Not In Use (Commented Out)!
%=====================================

% ACS format
% \isAPAandChicago{}{%
% \begin{thebibliography}{999}
% % Reference 1
% \bibitem[Author1(year)]{ref-journal}
% Author~1, T. The title of the cited article. {\em Journal Abbreviation} {\bf 2008}, {\em 10}, 142--149.
% % Reference 2
% \bibitem[Author2(year)]{ref-book1}
% Author~2, L. The title of the cited contribution. In {\em The Book Title}; Editor 1, F., Editor 2, A., Eds.; Publishing House: City, Country, 2007; pp. 32--58.
% % Reference 3
% \bibitem[Author1 and Author2 (year)]{ref-book2}
% Author 1, A.; Author 2, B. \textit{Book Title}, 3rd ed.; Publisher: Publisher Location, Country, 2008; pp. 154--196.
% % Reference 4
% \bibitem[Author4(year)]{ref-unpublish}
% Author 1, A.B.; Author 2, C. Title of Unpublished Work. \textit{Abbreviated Journal Name} year, \textit{phrase indicating stage of publication (submitted; accepted; in press)}.
% % Reference 5
% \bibitem[Author8(year)]{ref-url}
% Title of Site. Available online: URL (accessed on Day Month Year).
% % Reference 6
% \bibitem[Author6(year)]{ref-proceeding}
% Author 1, A.B.; Author 2, C.D.; Author 3, E.F. Title of presentation. In Proceedings of the Name of the Conference, Location of Conference, Country, Date of Conference (Day Month Year); Abstract Number (optional), Pagination (optional).
% % Reference 7
% \bibitem[Author7(year)]{ref-thesis}
% Author 1, A.B. Title of Thesis. Level of Thesis, Degree-Granting University, Location of University, Date of Completion.
% \end{thebibliography}
% }

% Chicago format (Used for journal: arts, genealogy, histories, humanities, jintelligence, laws, literature, religions, risks, socsci)
\isChicagoStyle{%
\begin{thebibliography}{999}
% Reference 1
\bibitem[Aranceta-Bartrina(1999a)]{ref-journal}
Aranceta-Bartrina, Javier. 1999a. Title of the cited article. \textit{Journal Title} 6: 100--10.
% Reference 2
\bibitem[Aranceta-Bartrina(1999b)]{ref-book1}
Aranceta-Bartrina, Javier. 1999b. Title of the chapter. In \textit{Book Title}, 2nd ed. Edited by Editor 1 and Editor 2. Publication place: Publisher, vol. 3, pp. 54–96.
% Reference 3
\bibitem[Baranwal and Munteanu {[1921]}(1955)]{ref-book2}
Baranwal, Ajay K., and Costea Munteanu. 1955. \textit{Book Title}. Publication place: Publisher, pp. 154--96. First published 1921 (op-tional).
% Reference 4
\bibitem[Berry and Smith(1999)]{ref-thesis}
Berry, Evan, and Amy M. Smith. 1999. Title of Thesis. Level of Thesis, Degree-Granting University, City, Country. Identifi-cation information (if available).
% Reference 5
\bibitem[Cojocaru et al.(1999)]{ref-unpublish}
Cojocaru, Ludmila, Dragos Constatin Sanda, and Eun Kyeong Yun. 1999. Title of Unpublished Work. \textit{Journal Title}, phrase indicating stage of publication.
% Reference 6
\bibitem[Driver et al.(2000)]{ref-proceeding}
Driver, John P., Steffen Rohrs, and Sean Meighoo. 2000. Title of Presentation. In \textit{Title of the Collected Work} (if available). Paper presented at Name of the Conference, Location of Conference, Date of Conference.
% Reference 7
\bibitem[Harwood(2008)]{ref-url}
Harwood, John. 2008. Title of the cited article. Available online: URL (accessed on Day Month Year).
\end{thebibliography}
}{}

% APA format (Used for journal: admsci, behavsci, businesses, econometrics, economies, education, ejihpe, games, humans, ijfs, journalmedia, jrfm, languages, psycholint, publications, tourismhosp, youth)
\isAPAStyle{%
\begin{thebibliography}{999}
% Reference 1
\bibitem[\protect\citeauthoryear{Azikiwe \BBA\ Bello}{{2020a}}]{ref-journal}
Azikiwe, H., \& Bello, A. (2020a). Title of the cited article. \textit{Journal Title}, \textit{Volume}(Issue), 
Firstpage--Lastpage/Article Number.
% Reference 2
\bibitem[\protect\citeauthoryear{Azikiwe \BBA\ Bello}{{2020b}}]{ref-book1}
Azikiwe, H., \& Bello, A. (2020b). \textit{Book title}. Publisher Name.
% Reference 3
\bibitem[Davison(1623/2019)]{ref-book2}
Davison, T. E. (2019). Title of the book chapter. In A. A. Editor (Ed.), \textit{Title of the book: Subtitle} 
(pp. Firstpage--Lastpage). Publisher Name. (Original work published 1623) (Optional).
% Reference 4
\bibitem[Fistek et al.(2017)]{ref-proceeding}
Fistek, A., Jester, E., \& Sonnenberg, K. (2017, Month Day). Title of contribution [Type of contribution]. Conference Name, Conference City, Conference Country.
% Reference 5
\bibitem[Hutcheson(2012)]{ref-thesis}
Hutcheson, V. H. (2012). \textit{Title of the thesis} [XX Thesis, Name of Institution Awarding the Degree].
% Reference 6
\bibitem[Lippincott \& Poindexter(2019)]{ref-unpublish}
Lippincott, T., \& Poindexter, E. K. (2019). \textit{Title of the unpublished manuscript} [Unpublished manuscript/Manuscript in prepara-tion/Manuscript submitted for publication]. Department Name, Institution Name.
% Reference 7
\bibitem[Harwood(2008)]{ref-url}
Harwood, J. (2008). \textit{Title of the cited article}. Available online: URL (accessed on Day Month Year).
\end{thebibliography}
}{}

% If authors have biography, please use the format below
%\section*{Short Biography of Authors}
%\bio
%{\raisebox{-0.35cm}{\includegraphics[width=3.5cm,height=5.3cm,clip,keepaspectratio]{Definitions/author1.pdf}}}
%{\textbf{Firstname Lastname} Biography of first author}
%
%\bio
%{\raisebox{-0.35cm}{\includegraphics[width=3.5cm,height=5.3cm,clip,keepaspectratio]{Definitions/author2.jpg}}}
%{\textbf{Firstname Lastname} Biography of second author}

% For the MDPI journals use author-date citation, please follow the formatting guidelines on http://www.mdpi.com/authors/references
% To cite two works by the same author: \citeauthor{ref-journal-1a} (\citeyear{ref-journal-1a}, \citeyear{ref-journal-1b}). This produces: Whittaker (1967, 1975)
% To cite two works by the same author with specific pages: \citeauthor{ref-journal-3a} (\citeyear{ref-journal-3a}, p. 328; \citeyear{ref-journal-3b}, p.475). This produces: Wong (1999, p. 328; 2000, p. 475)

%%%%%%%%%%%%%%%%%%%%%%%%%%%%%%%%%%%%%%%%%%
%% for journal Sci
%\reviewreports{\\
%Reviewer 1 comments and authors’ response\\
%Reviewer 2 comments and authors’ response\\
%Reviewer 3 comments and authors’ response
%}
%%%%%%%%%%%%%%%%%%%%%%%%%%%%%%%%%%%%%%%%%%
\PublishersNote{}
%\isPreprints{}{% This command is only used for ``preprints''.
\end{adjustwidth}
%} % If the paper is ``preprints'', please uncomment this parenthesis.
\end{document}

